\documentclass[12pt]{article}
\usepackage[a4paper,margin=1in]{geometry}
\usepackage{iftex}
\ifXeTeX
\usepackage{fontspec}
\defaultfontfeatures{Ligatures=TeX}
\setmainfont{Times New Roman}
\else
\usepackage[T1]{fontenc}
\usepackage[utf8]{inputenc}
\usepackage{newtxtext}
\usepackage{newtxmath}
\fi
\usepackage{enumitem}
\usepackage{hyperref}
\usepackage{parskip}
\usepackage{microtype}
\usepackage[normalem]{ulem}
\usepackage{xcolor}

\newcommand{\practiceid}[1]{\texttt{#1}}
\newlength{\readinggutter}
\setlength{\readinggutter}{2.8em}
\newlength{\readingfirstindent}
\setlength{\readingfirstindent}{1.5em}
\newlength{\questionblockgap}
\setlength{\questionblockgap}{0.45\baselineskip}
\newlength{\exampleblockgap}
\setlength{\exampleblockgap}{0.65\baselineskip}
\newlength{\passageparagraphgap}
\setlength{\passageparagraphgap}{0.45\baselineskip}
\newenvironment{exampleblock}{\par\vspace{0.35\baselineskip}\begingroup\raggedright\setlength{\parskip}{0pt}\setlength{\parindent}{0pt}}{\par\endgroup\vspace{\exampleblockgap}}
\newcommand{\readinglineindent}[2]{%
\par\noindent%
\hangindent=\dimexpr\readinggutter+0.9em\relax%
\hangafter=1%
\makebox[\readinggutter][r]{\scriptsize\color{gray}#1}\hspace{0.9em}\hspace*{\readingfirstindent}#2%
\par}
\newcommand{\readingline}[2]{%
\par\noindent%
\hangindent=\dimexpr\readinggutter+0.9em\relax%
\hangafter=1%
\makebox[\readinggutter][r]{\scriptsize\color{gray}#1}\hspace{0.9em}#2%
\par}

\setlength{\parindent}{0pt}
\setlength{\parskip}{0.5em}
\setlist[enumerate]{topsep=0.25em,itemsep=0.18em,parsep=0pt,partopsep=0pt}
\setlength{\emergencystretch}{2em}

\begin{document}
\section*{TOEFL Practice 04 - Review}
Source of truth: DOCX only.
\section*{Listening Comprehension}
\subsection*{Part A}
DIRECTIONS: In Part A you will hear short conversations between two speakers. At the end of each conversation, a third speaker will ask a question about what the first two speakers said. Each conversation and each question will be spoken only one time. Therefore, you must listen carefully to understand what each speaker says. After you hear a conversation and the question, read the four choices and select the one that is the best answer to the question the speaker asked. Then, on your answer. sheet, find the number of the question and blacken the space that corresponds to the letter for the answer you have chosen. Blacken the space completely so that the letter inside the space does not show.\par
\begin{exampleblock}
Listen to the following example.\par
On the recording, you hear:\par
(Man) Does the car need to be filled?\par
(Woman) Mary stopped at the gas station on her way home.\par
(Narrator) What does the woman mean?\par
In your test book, you will read:\par
(A) Mary bought some food.\par
(B) Mary had car trouble.\par
(C) Mary went shopping.\par
(D) Mary bought some gas.\par
From the conversation you learn that Mary stopped at the gas station on her way home. The best answer to the question "Does the car need to be filled?" is (D), "Mary bought some gas." Therefore, the correct answer is (D).\par
Now let us begin Part A with question number 1.\par
\end{exampleblock}
\begin{exampleblock}
\textbf{Transkrip rekonstruksi AI (draft; bukan resmi)}\par
Man: I heard Professor White's advanced seminar is one of the most demanding courses in the department.\par
Woman: It is, but his students always speak of him with real admiration.\par
Narrator: What does the woman imply?\par
\end{exampleblock}
\noindent\textbf{1.}
\begin{enumerate}[label=(\Alph*),leftmargin=*,itemsep=0.2em]
\item Professor White holds advanced classes.
\item Professor White's students have graduated.
\item Students have great respect for Professor White.
\item Students are taking Professor White's seminar.
\end{enumerate}
\par
\begingroup
\setlength{\parskip}{0pt}
\setlength{\parindent}{0pt}
\textbf{Review}\par
Answer: C\par
Confidence: medium\par
Jawaban diambil dari kunci belajar yang diberikan pengguna dan belum diverifikasi terhadap audio/transkrip resmi.\par
\endgroup
\par
\vspace{0.25\baselineskip}
\begin{exampleblock}
\textbf{Transkrip rekonstruksi AI (draft; bukan resmi)}\par
Man: That was an excellent dinner.\par
Woman: I'm glad you enjoyed it. Would you care for some dessert?\par
Narrator: What does the woman mean?\par
\end{exampleblock}
\noindent\textbf{2.}
\begin{enumerate}[label=(\Alph*),leftmargin=*,itemsep=0.2em]
\item You have to be careful in the desert.
\item You shouldn't eat so many sweets.
\item May I offer you some dessert?
\item Can you help me with the sweets?
\end{enumerate}
\par
\begingroup
\setlength{\parskip}{0pt}
\setlength{\parindent}{0pt}
\textbf{Review}\par
Answer: C\par
Confidence: medium\par
Jawaban diambil dari kunci belajar yang diberikan pengguna dan belum diverifikasi terhadap audio/transkrip resmi.\par
\endgroup
\par
\vspace{0.25\baselineskip}
\begin{exampleblock}
\textbf{Transkrip rekonstruksi AI (draft; bukan resmi)}\par
Woman: Could you lend me a little money for the vending machine?\par
Man: Sorry, I left my wallet at home.\par
Narrator: What does the man mean?\par
\end{exampleblock}
\noindent\textbf{3.}
\begin{enumerate}[label=(\Alph*),leftmargin=*,itemsep=0.2em]
\item He never has any money.
\item He has no money on him.
\item He is out getting cash.
\item He has been out of work.
\end{enumerate}
\par
\begingroup
\setlength{\parskip}{0pt}
\setlength{\parindent}{0pt}
\textbf{Review}\par
Answer: B\par
Confidence: medium\par
Jawaban diambil dari kunci belajar yang diberikan pengguna dan belum diverifikasi terhadap audio/transkrip resmi.\par
\endgroup
\par
\vspace{0.25\baselineskip}
\begin{exampleblock}
\textbf{Transkrip rekonstruksi AI (draft; bukan resmi)}\par
Man: Should I leave the copying machine on after I finish these copies?\par
Woman: No, please turn it off when you're done.\par
Narrator: What does the woman mean?\par
\end{exampleblock}
\noindent\textbf{4.}
\begin{enumerate}[label=(\Alph*),leftmargin=*,itemsep=0.2em]
\item The machine shouldn't be left running.
\item Turn on the copying machine when you leave.
\item The machine has to be turned around.
\item You made enough copies already.
\end{enumerate}
\par
\begingroup
\setlength{\parskip}{0pt}
\setlength{\parindent}{0pt}
\textbf{Review}\par
Answer: A\par
Confidence: medium\par
Jawaban diambil dari kunci belajar yang diberikan pengguna dan belum diverifikasi terhadap audio/transkrip resmi.\par
\endgroup
\par
\vspace{0.25\baselineskip}
\begin{exampleblock}
\textbf{Transkrip rekonstruksi AI (draft; bukan resmi)}\par
Man: Did you hear that Tom passed his Latin course?\par
Woman: Really? I thought everyone expected him to fail it.\par
Narrator: What does the woman imply?\par
\end{exampleblock}
\noindent\textbf{5.}
\begin{enumerate}[label=(\Alph*),leftmargin=*,itemsep=0.2em]
\item Tom passed up a Latin course.
\item Tom was expected to fail Latin.
\item He's surprised Tom is late.
\item Tom's Latin class is last.
\end{enumerate}
\par
\begingroup
\setlength{\parskip}{0pt}
\setlength{\parindent}{0pt}
\textbf{Review}\par
Answer: B\par
Confidence: medium\par
Jawaban diambil dari kunci belajar yang diberikan pengguna dan belum diverifikasi terhadap audio/transkrip resmi.\par
\endgroup
\par
\vspace{0.25\baselineskip}
\begin{exampleblock}
\textbf{Transkrip rekonstruksi AI (draft; bukan resmi)}\par
Man: Did she spend one hundred fifteen dollars on that dress?\par
Woman: The dress was one hundred fifteen, and the matching jacket was ninety-five.\par
Narrator: How much did she spend?\par
\end{exampleblock}
\noindent\textbf{6.}
\begin{enumerate}[label=(\Alph*),leftmargin=*,itemsep=0.2em]
\item She buys exotic clothes.
\item She bought a \$115 dress.
\item She spent \$210.
\item She has a new house.
\end{enumerate}
\par
\begingroup
\setlength{\parskip}{0pt}
\setlength{\parindent}{0pt}
\textbf{Review}\par
Answer: C\par
Confidence: medium\par
Jawaban diambil dari kunci belajar yang diberikan pengguna dan belum diverifikasi terhadap audio/transkrip resmi.\par
\endgroup
\par
\vspace{0.25\baselineskip}
\begin{exampleblock}
\textbf{Transkrip rekonstruksi AI (draft; bukan resmi)}\par
Man: What do you think of the new engineering building?\par
Woman: To tell the truth, it isn't at all what I expected.\par
Narrator: What does the woman mean?\par
\end{exampleblock}
\noindent\textbf{7.}
\begin{enumerate}[label=(\Alph*),leftmargin=*,itemsep=0.2em]
\item The new engineering building is far from here.
\item The building is not what I thought it would be.
\item I don't know what I expected the engineer to do.
\item They are building a new house for the engineer.
\end{enumerate}
\par
\begingroup
\setlength{\parskip}{0pt}
\setlength{\parindent}{0pt}
\textbf{Review}\par
Answer: B\par
Confidence: medium\par
Jawaban diambil dari kunci belajar yang diberikan pengguna dan belum diverifikasi terhadap audio/transkrip resmi.\par
\endgroup
\par
\vspace{0.25\baselineskip}
\begin{exampleblock}
\textbf{Transkrip rekonstruksi AI (draft; bukan resmi)}\par
Woman: Have you finally made that appointment with the dentist?\par
Man: Not yet. I keep putting it off.\par
Narrator: What does the man mean?\par
\end{exampleblock}
\noindent\textbf{8.}
\begin{enumerate}[label=(\Alph*),leftmargin=*,itemsep=0.2em]
\item The man has put off going to a dentist.
\item The man likes getting his teeth cleaned.
\item She is not happy that the man went to a dentist.
\item Everyone should have their teeth cleaned regularly.
\end{enumerate}
\par
\begingroup
\setlength{\parskip}{0pt}
\setlength{\parindent}{0pt}
\textbf{Review}\par
Answer: A\par
Confidence: medium\par
Jawaban diambil dari kunci belajar yang diberikan pengguna dan belum diverifikasi terhadap audio/transkrip resmi.\par
\endgroup
\par
\vspace{0.25\baselineskip}
\begin{exampleblock}
\textbf{Transkrip rekonstruksi AI (draft; bukan resmi)}\par
Man: Do you think this hall will work for the wedding reception?\par
Woman: Not for two hundred guests. It won't hold nearly that many people.\par
Narrator: What does the woman imply?\par
\end{exampleblock}
\noindent\textbf{9.}
\begin{enumerate}[label=(\Alph*),leftmargin=*,itemsep=0.2em]
\item The rent for the hall is high.
\item The wedding hall is for rent.
\item The hall is too small.
\item There are not enough guests.
\end{enumerate}
\par
\begingroup
\setlength{\parskip}{0pt}
\setlength{\parindent}{0pt}
\textbf{Review}\par
Answer: C\par
Confidence: medium\par
Jawaban diambil dari kunci belajar yang diberikan pengguna dan belum diverifikasi terhadap audio/transkrip resmi.\par
\endgroup
\par
\vspace{0.25\baselineskip}
\begin{exampleblock}
\textbf{Transkrip rekonstruksi AI (draft; bukan resmi)}\par
Man: Could you cash this traveler's check for one thousand dollars?\par
Woman: I'm sorry. For checks that large, only the manager is authorized to cash them.\par
Narrator: What does the woman mean?\par
\end{exampleblock}
\noindent\textbf{10.}
\begin{enumerate}[label=(\Alph*),leftmargin=*,itemsep=0.2em]
\item She can cash the check.
\item Only the manager can cash the check.
\item Traveller's checks cannot be cashed there.
\item A thousand dollars is a large sum of money.
\end{enumerate}
\par
\begingroup
\setlength{\parskip}{0pt}
\setlength{\parindent}{0pt}
\textbf{Review}\par
Answer: B\par
Confidence: medium\par
Jawaban diambil dari kunci belajar yang diberikan pengguna dan belum diverifikasi terhadap audio/transkrip resmi.\par
\endgroup
\par
\vspace{0.25\baselineskip}
\begin{exampleblock}
\textbf{Transkrip rekonstruksi AI (draft; bukan resmi)}\par
Man: How did the housecleaning go yesterday?\par
Woman: Great. The crew washed all the windows and cleaned the rugs.\par
Narrator: What does the woman mean?\par
\end{exampleblock}
\noindent\textbf{11.}
\begin{enumerate}[label=(\Alph*),leftmargin=*,itemsep=0.2em]
\item They worked together as a team.
\item The windows and carpets are beautiful.
\item They now have clean windows and rugs.
\item They are always cleaning their house.
\end{enumerate}
\par
\begingroup
\setlength{\parskip}{0pt}
\setlength{\parindent}{0pt}
\textbf{Review}\par
Answer: C\par
Confidence: medium\par
Jawaban diambil dari kunci belajar yang diberikan pengguna dan belum diverifikasi terhadap audio/transkrip resmi.\par
\endgroup
\par
\vspace{0.25\baselineskip}
\begin{exampleblock}
\textbf{Transkrip rekonstruksi AI (draft; bukan resmi)}\par
Man: What was Linda's reaction when she won the award?\par
Woman: She was completely at a loss for words.\par
Narrator: What does the woman mean?\par
\end{exampleblock}
\noindent\textbf{12.}
\begin{enumerate}[label=(\Alph*),leftmargin=*,itemsep=0.2em]
\item She got lost one more time.
\item She didn't know what to say.
\item She is learning many new words.
\item She had less than the price of the book.
\end{enumerate}
\par
\begingroup
\setlength{\parskip}{0pt}
\setlength{\parindent}{0pt}
\textbf{Review}\par
Answer: B\par
Confidence: medium\par
Jawaban diambil dari kunci belajar yang diberikan pengguna dan belum diverifikasi terhadap audio/transkrip resmi.\par
\endgroup
\par
\vspace{0.25\baselineskip}
\begin{exampleblock}
\textbf{Transkrip rekonstruksi AI (draft; bukan resmi)}\par
Woman: Mark looks exhausted this morning.\par
Man: He really should stop keeping such late hours.\par
Narrator: What does the man mean about Mark?\par
\end{exampleblock}
\noindent\textbf{13.}
\begin{enumerate}[label=(\Alph*),leftmargin=*,itemsep=0.2em]
\item He is staying in this hotel.
\item He keeps going to bed late.
\item He lives by himself in this house.
\item He started early and finished late.
\end{enumerate}
\par
\begingroup
\setlength{\parskip}{0pt}
\setlength{\parindent}{0pt}
\textbf{Review}\par
Answer: B\par
Confidence: medium\par
Jawaban diambil dari kunci belajar yang diberikan pengguna dan belum diverifikasi terhadap audio/transkrip resmi.\par
\endgroup
\par
\vspace{0.25\baselineskip}
\begin{exampleblock}
\textbf{Transkrip rekonstruksi AI (draft; bukan resmi)}\par
Man: Is Ronald planning to travel this summer?\par
Woman: Yes, he says he wants to go abroad, but he hasn't decided which country yet.\par
Narrator: What does the woman mean?\par
\end{exampleblock}
\noindent\textbf{14.}
\begin{enumerate}[label=(\Alph*),leftmargin=*,itemsep=0.2em]
\item Ronald won't travel this summer.
\item She doesn't know to what country Ronald is going.
\item She wants to travel with Ronald.
\item Ronald wants to travel out of the country.
\end{enumerate}
\par
\begingroup
\setlength{\parskip}{0pt}
\setlength{\parindent}{0pt}
\textbf{Review}\par
Answer: D\par
Confidence: medium\par
Jawaban diambil dari kunci belajar yang diberikan pengguna dan belum diverifikasi terhadap audio/transkrip resmi.\par
\endgroup
\par
\vspace{0.25\baselineskip}
\begin{exampleblock}
\textbf{Transkrip rekonstruksi AI (draft; bukan resmi)}\par
Woman: The store said they couldn't deliver my dress today.\par
Man: Did you give them your address? They can't deliver it if they don't know where you live.\par
Narrator: What does the man imply?\par
\end{exampleblock}
\noindent\textbf{15.}
\begin{enumerate}[label=(\Alph*),leftmargin=*,itemsep=0.2em]
\item They can't deliver her dress.
\item They don't know where she lives.
\item They didn't touch her dress.
\item They aren't familiar with the city.
\end{enumerate}
\par
\begingroup
\setlength{\parskip}{0pt}
\setlength{\parindent}{0pt}
\textbf{Review}\par
Answer: B\par
Confidence: medium\par
Jawaban diambil dari kunci belajar yang diberikan pengguna dan belum diverifikasi terhadap audio/transkrip resmi.\par
\endgroup
\par
\vspace{0.25\baselineskip}
\begin{exampleblock}
\textbf{Transkrip rekonstruksi AI (draft; bukan resmi)}\par
Man: We still have forty miles to go, and the meeting starts in half an hour.\par
Woman: Then there is no way we'll get there on time.\par
Narrator: What does the woman mean?\par
\end{exampleblock}
\noindent\textbf{16.}
\begin{enumerate}[label=(\Alph*),leftmargin=*,itemsep=0.2em]
\item We'll be there on time.
\item We'll be late.
\item We'll take the next turn.
\item The turn is forty miles away.
\end{enumerate}
\par
\begingroup
\setlength{\parskip}{0pt}
\setlength{\parindent}{0pt}
\textbf{Review}\par
Answer: B\par
Confidence: medium\par
Jawaban diambil dari kunci belajar yang diberikan pengguna dan belum diverifikasi terhadap audio/transkrip resmi.\par
\endgroup
\par
\vspace{0.25\baselineskip}
\begin{exampleblock}
\textbf{Transkrip rekonstruksi AI (draft; bukan resmi)}\par
Man: The business school expected nine hundred new students this year.\par
Woman: But only about three hundred actually enrolled.\par
Narrator: What does the woman imply?\par
\end{exampleblock}
\noindent\textbf{17.}
\begin{enumerate}[label=(\Alph*),leftmargin=*,itemsep=0.2em]
\item Only a third of the students are enrolled in the business school.
\item The third enrollment prediction for the business school was accurate.
\item The enrollment is three times higher than what was predicted.
\item The enrollment figures are about 66 percent lower than expected.
\end{enumerate}
\par
\begingroup
\setlength{\parskip}{0pt}
\setlength{\parindent}{0pt}
\textbf{Review}\par
Answer: D\par
Confidence: medium\par
Jawaban diambil dari kunci belajar yang diberikan pengguna dan belum diverifikasi terhadap audio/transkrip resmi.\par
\endgroup
\par
\vspace{0.25\baselineskip}
\begin{exampleblock}
\textbf{Transkrip rekonstruksi AI (draft; bukan resmi)}\par
Man: Did you hear that Mr. Benson has started taking singing lessons?\par
Woman: At his age? That's certainly an unusual time to begin.\par
Narrator: What does the woman mean?\par
\end{exampleblock}
\noindent\textbf{18.}
\begin{enumerate}[label=(\Alph*),leftmargin=*,itemsep=0.2em]
\item He likes to sing old songs.
\item The rumors about him are unbelievable.
\item To start singing at his age is unusual.
\item You shouldn't talk about him this way.
\end{enumerate}
\par
\begingroup
\setlength{\parskip}{0pt}
\setlength{\parindent}{0pt}
\textbf{Review}\par
Answer: C\par
Confidence: medium\par
Jawaban diambil dari kunci belajar yang diberikan pengguna dan belum diverifikasi terhadap audio/transkrip resmi.\par
\endgroup
\par
\vspace{0.25\baselineskip}
\begin{exampleblock}
\textbf{Transkrip rekonstruksi AI (draft; bukan resmi)}\par
Man: I explained the assignment three times, but the students still asked the same questions afterward.\par
Woman: It sounds as if they aren't paying attention to what you say.\par
Narrator: What does the woman mean?\par
\end{exampleblock}
\noindent\textbf{19.}
\begin{enumerate}[label=(\Alph*),leftmargin=*,itemsep=0.2em]
\item He doesn't feel good, and he won't go to class.
\item He couldn't hear the lecture because students were talking.
\item Students don't pay attention to what he says.
\item In this class, the same thing happens every day.
\end{enumerate}
\par
\begingroup
\setlength{\parskip}{0pt}
\setlength{\parindent}{0pt}
\textbf{Review}\par
Answer: C\par
Confidence: medium\par
Jawaban diambil dari kunci belajar yang diberikan pengguna dan belum diverifikasi terhadap audio/transkrip resmi.\par
\endgroup
\par
\vspace{0.25\baselineskip}
\begin{exampleblock}
\textbf{Transkrip rekonstruksi AI (draft; bukan resmi)}\par
Man: What did you and your sister do after dinner?\par
Woman: We took a walk along the river and watched the sunset.\par
Narrator: What did the woman and her sister do?\par
\end{exampleblock}
\noindent\textbf{20.}
\begin{enumerate}[label=(\Alph*),leftmargin=*,itemsep=0.2em]
\item They passed the river.
\item They went for a walk.
\item They visited her sister.
\item They were looking at the sunset.
\end{enumerate}
\par
\begingroup
\setlength{\parskip}{0pt}
\setlength{\parindent}{0pt}
\textbf{Review}\par
Answer: B\par
Confidence: medium\par
Jawaban diambil dari kunci belajar yang diberikan pengguna dan belum diverifikasi terhadap audio/transkrip resmi.\par
\endgroup
\par
\vspace{0.25\baselineskip}
\begin{exampleblock}
\textbf{Transkrip rekonstruksi AI (draft; bukan resmi)}\par
Man: Should we go down to the office for tickets to the game?\par
Woman: Let's call first and see whether any are still available.\par
Narrator: What will they probably do first?\par
\end{exampleblock}
\noindent\textbf{21.}
\begin{enumerate}[label=(\Alph*),leftmargin=*,itemsep=0.2em]
\item Stand in line
\item Try to order tickets
\item Go to the game
\item Call their office
\end{enumerate}
\par
\begingroup
\setlength{\parskip}{0pt}
\setlength{\parindent}{0pt}
\textbf{Review}\par
Answer: B\par
Confidence: medium\par
Jawaban diambil dari kunci belajar yang diberikan pengguna dan belum diverifikasi terhadap audio/transkrip resmi.\par
\endgroup
\par
\vspace{0.25\baselineskip}
\begin{exampleblock}
\textbf{Transkrip rekonstruksi AI (draft; bukan resmi)}\par
Man: Since the phone is always busy, maybe we should get another line.\par
Woman: I don't think we can afford another monthly bill right now.\par
Narrator: What does the woman mean?\par
\end{exampleblock}
\noindent\textbf{22.}
\begin{enumerate}[label=(\Alph*),leftmargin=*,itemsep=0.2em]
\item She doesn't need another phone.
\item She is not as busy as it seems.
\item They cant afford another line.
\item They have to use the phone less.
\end{enumerate}
\par
\begingroup
\setlength{\parskip}{0pt}
\setlength{\parindent}{0pt}
\textbf{Review}\par
Answer: C\par
Confidence: medium\par
Jawaban diambil dari kunci belajar yang diberikan pengguna dan belum diverifikasi terhadap audio/transkrip resmi.\par
\endgroup
\par
\vspace{0.25\baselineskip}
\begin{exampleblock}
\textbf{Transkrip rekonstruksi AI (draft; bukan resmi)}\par
Woman: Since I already took math at another school, can I skip the placement test?\par
Man: No. All new students have to take the math placement test, even with transfer credit.\par
Narrator: What does the man mean?\par
\end{exampleblock}
\noindent\textbf{23.}
\begin{enumerate}[label=(\Alph*),leftmargin=*,itemsep=0.2em]
\item The man administers math placement tests.
\item The woman should take a math course.
\item The woman has to take the placement test.
\item The woman can transfer her math credit.
\end{enumerate}
\par
\begingroup
\setlength{\parskip}{0pt}
\setlength{\parindent}{0pt}
\textbf{Review}\par
Answer: C\par
Confidence: medium\par
Jawaban diambil dari kunci belajar yang diberikan pengguna dan belum diverifikasi terhadap audio/transkrip resmi.\par
\endgroup
\par
\vspace{0.25\baselineskip}
\begin{exampleblock}
\textbf{Transkrip rekonstruksi AI (draft; bukan resmi)}\par
Man: The roses need pruning, and the lawn has to be reseeded.\par
Woman: Then we should hire someone who really knows how to care for plants.\par
Narrator: Whom do they need?\par
\end{exampleblock}
\noindent\textbf{24.}
\begin{enumerate}[label=(\Alph*),leftmargin=*,itemsep=0.2em]
\item A florist
\item A gardener
\item A barber
\item A custodian
\end{enumerate}
\par
\begingroup
\setlength{\parskip}{0pt}
\setlength{\parindent}{0pt}
\textbf{Review}\par
Answer: B\par
Confidence: medium\par
Jawaban diambil dari kunci belajar yang diberikan pengguna dan belum diverifikasi terhadap audio/transkrip resmi.\par
\endgroup
\par
\vspace{0.25\baselineskip}
\begin{exampleblock}
\textbf{Transkrip rekonstruksi AI (draft; bukan resmi)}\par
Man: What are you doing in the garage?\par
Woman: I'm just listening to the final round of the contest on the radio.\par
Narrator: What is the woman doing?\par
\end{exampleblock}
\noindent\textbf{25.}
\begin{enumerate}[label=(\Alph*),leftmargin=*,itemsep=0.2em]
\item Listening to the radio
\item Watching a contest
\item Repairing a car battery
\item Attending a conference
\end{enumerate}
\par
\begingroup
\setlength{\parskip}{0pt}
\setlength{\parindent}{0pt}
\textbf{Review}\par
Answer: A\par
Confidence: medium\par
Jawaban diambil dari kunci belajar yang diberikan pengguna dan belum diverifikasi terhadap audio/transkrip resmi.\par
\endgroup
\par
\vspace{0.25\baselineskip}
\begin{exampleblock}
\textbf{Transkrip rekonstruksi AI (draft; bukan resmi)}\par
Man: I'd like to rent a compact car for the weekend.\par
Woman: Certainly. May I see your driver's license and credit card?\par
Narrator: Where does this conversation probably take place?\par
\end{exampleblock}
\noindent\textbf{26.}
\begin{enumerate}[label=(\Alph*),leftmargin=*,itemsep=0.2em]
\item An insurance company
\item A car rental agency
\item A real estate agency
\item An apartment complex
\end{enumerate}
\par
\begingroup
\setlength{\parskip}{0pt}
\setlength{\parindent}{0pt}
\textbf{Review}\par
Answer: B\par
Confidence: medium\par
Jawaban diambil dari kunci belajar yang diberikan pengguna dan belum diverifikasi terhadap audio/transkrip resmi.\par
\endgroup
\par
\vspace{0.25\baselineskip}
\begin{exampleblock}
\textbf{Transkrip rekonstruksi AI (draft; bukan resmi)}\par
Man: I feel sick, but my exam is today. Do you think I should stay home or go take it?\par
Woman: I wish I could tell you, but I really don't know what you should do.\par
Narrator: What does the woman mean?\par
\end{exampleblock}
\noindent\textbf{27.}
\begin{enumerate}[label=(\Alph*),leftmargin=*,itemsep=0.2em]
\item The man shouldn't attend the exam.
\item The instructor isn't proctoring the exam.
\item In this situation, the man should stay home.
\item She doesn't know what the man should do.
\end{enumerate}
\par
\begingroup
\setlength{\parskip}{0pt}
\setlength{\parindent}{0pt}
\textbf{Review}\par
Answer: D\par
Confidence: medium\par
Jawaban diambil dari kunci belajar yang diberikan pengguna dan belum diverifikasi terhadap audio/transkrip resmi.\par
\endgroup
\par
\vspace{0.25\baselineskip}
\begin{exampleblock}
\textbf{Transkrip rekonstruksi AI (draft; bukan resmi)}\par
Man: This course is impossible. I blame the class for my poor grade.\par
Woman: I don't think the course is at fault. It's difficult, but it is fair.\par
Narrator: What does the woman mean?\par
\end{exampleblock}
\noindent\textbf{28.}
\begin{enumerate}[label=(\Alph*),leftmargin=*,itemsep=0.2em]
\item She is a demanding person.
\item She is sympathetic to the man.
\item She doesn't fault the course.
\item She got out of the course.
\end{enumerate}
\par
\begingroup
\setlength{\parskip}{0pt}
\setlength{\parindent}{0pt}
\textbf{Review}\par
Answer: C\par
Confidence: medium\par
Jawaban diambil dari kunci belajar yang diberikan pengguna dan belum diverifikasi terhadap audio/transkrip resmi.\par
\endgroup
\par
\vspace{0.25\baselineskip}
\begin{exampleblock}
\textbf{Transkrip rekonstruksi AI (draft; bukan resmi)}\par
Man: Do you have the new edition of the biology textbook?\par
Woman: It's on the shelf with the other required course books. The used copies are in the next aisle.\par
Narrator: Where does this conversation probably take place?\par
\end{exampleblock}
\noindent\textbf{29.}
\begin{enumerate}[label=(\Alph*),leftmargin=*,itemsep=0.2em]
\item A bookstore
\item A publishing house
\item A craft show
\item An art exhibition
\end{enumerate}
\par
\begingroup
\setlength{\parskip}{0pt}
\setlength{\parindent}{0pt}
\textbf{Review}\par
Answer: A\par
Confidence: medium\par
Jawaban diambil dari kunci belajar yang diberikan pengguna dan belum diverifikasi terhadap audio/transkrip resmi.\par
\endgroup
\par
\vspace{0.25\baselineskip}
\begin{exampleblock}
\textbf{Transkrip rekonstruksi AI (draft; bukan resmi)}\par
Woman: Congratulations! I heard you came in first yesterday.\par
Man: Thanks. I still can't believe I won the cross-country ski race.\par
Narrator: What did the man do?\par
\end{exampleblock}
\noindent\textbf{30.}
\begin{enumerate}[label=(\Alph*),leftmargin=*,itemsep=0.2em]
\item He walked ten miles.
\item He came here first.
\item He won a ski competition.
\item He lives across the country.
\end{enumerate}
\par
\begingroup
\setlength{\parskip}{0pt}
\setlength{\parindent}{0pt}
\textbf{Review}\par
Answer: C\par
Confidence: medium\par
Jawaban diambil dari kunci belajar yang diberikan pengguna dan belum diverifikasi terhadap audio/transkrip resmi.\par
\endgroup
\par
\vspace{0.25\baselineskip}
\subsection*{Part B}
DIRECTIONS: In this part of the test, you will hear longer conversations. After each conversation, you will hear several questions. The conversations and questions will not be repeated. After you hear a question, read the four possible answers in. your test book and select. The best answer. Then, on your answer sheet, find the number of the question and fill in the space that corresponds to the letter of the answer you have chosen. Remember, you are not allowed to take notes or write in your test book.\par
\begin{exampleblock}
Listen to the following example:\par
You will hear:\par
You will read:\par
(A) He has changed jobs.\par
(B) He has two children.\par
(C) He has two jobs.\par
(D) He is looking for a job.\par
From the conversation you learn that Tom has taken an additional job. The best answer to the question "Why is Tom tired?" is (C), "He has two jobs." Therefore the correct answer is (C).\par
\end{exampleblock}
\begin{exampleblock}
\textbf{Transkrip rekonstruksi AI (draft; bukan resmi)}\par
Woman: The movie is almost ready. I'll make some popcorn before we start.\par
Man: Great idea. It's much better to have it here at home than to buy it at a concession stand.\par
Woman: Exactly. Could you get the large cooking pot from the cabinet? We don't need corn on the cob or any special heater, just a pot, some kernels, and a little oil.\par
Man: Popcorn really is one of the cheapest snacks around, isn't it?\par
Woman: It is. That's one reason people like it so much. But I won't add much butter. Butter tastes good, but it mostly adds extra calories.\par
Man: I've always wondered what actually makes popcorn pop.\par
Woman: Each kernel has a tiny amount of moisture inside. When the kernel is heated, that moisture turns to steam. Pressure builds until the kernel explodes.\par
Man: So what is the oil for? Does it make the corn pop?\par
Woman: Not exactly. The oil helps coat the kernels and transfer heat evenly. Without oil, the kernels could burn before they pop.\par
Man: And after they pop, they are huge compared with the original kernels.\par
Woman: Right. A popped kernel is many times larger than it was before.\par
\end{exampleblock}
\noindent\textbf{31.}
\begin{enumerate}[label=(\Alph*),leftmargin=*,itemsep=0.2em]
\item In a movie theater
\item In a baseball stadium
\item At the man and woman's house
\item At a concession stand
\end{enumerate}
\par
\begingroup
\setlength{\parskip}{0pt}
\setlength{\parindent}{0pt}
\textbf{Review}\par
Answer: C\par
Confidence: medium\par
Jawaban diambil dari kunci belajar yang diberikan pengguna dan belum diverifikasi terhadap audio/transkrip resmi.\par
\endgroup
\par
\vspace{0.25\baselineskip}
\noindent\textbf{32.}
\begin{enumerate}[label=(\Alph*),leftmargin=*,itemsep=0.2em]
\item Corn on the cob
\item Four types of oil
\item A cooking pot
\item A small heater
\end{enumerate}
\par
\begingroup
\setlength{\parskip}{0pt}
\setlength{\parindent}{0pt}
\textbf{Review}\par
Answer: C\par
Confidence: medium\par
Jawaban diambil dari kunci belajar yang diberikan pengguna dan belum diverifikasi terhadap audio/transkrip resmi.\par
\endgroup
\par
\vspace{0.25\baselineskip}
\noindent\textbf{33.}
\begin{enumerate}[label=(\Alph*),leftmargin=*,itemsep=0.2em]
\item They are not good for you.
\item They are cheap.
\item They were invented by settlers.
\item They need moisture.
\end{enumerate}
\par
\begingroup
\setlength{\parskip}{0pt}
\setlength{\parindent}{0pt}
\textbf{Review}\par
Answer: B\par
Confidence: medium\par
Jawaban diambil dari kunci belajar yang diberikan pengguna dan belum diverifikasi terhadap audio/transkrip resmi.\par
\endgroup
\par
\vspace{0.25\baselineskip}
\noindent\textbf{34.}
\begin{enumerate}[label=(\Alph*),leftmargin=*,itemsep=0.2em]
\item It's necessary.
\item It adds calories.
\item It's good for you.
\item It tastes terrible.
\end{enumerate}
\par
\begingroup
\setlength{\parskip}{0pt}
\setlength{\parindent}{0pt}
\textbf{Review}\par
Answer: B\par
Confidence: medium\par
Jawaban diambil dari kunci belajar yang diberikan pengguna dan belum diverifikasi terhadap audio/transkrip resmi.\par
\endgroup
\par
\vspace{0.25\baselineskip}
\noindent\textbf{35.}
\begin{enumerate}[label=(\Alph*),leftmargin=*,itemsep=0.2em]
\item The liquid coats the corn.
\item Oil puffs up the corn.
\item It makes the kernel explode.
\item Moisture expands at the seams.
\end{enumerate}
\par
\begingroup
\setlength{\parskip}{0pt}
\setlength{\parindent}{0pt}
\textbf{Review}\par
Answer: C\par
Confidence: medium\par
Jawaban diambil dari kunci belajar yang diberikan pengguna dan belum diverifikasi terhadap audio/transkrip resmi.\par
\endgroup
\par
\vspace{0.25\baselineskip}
\noindent\textbf{36.}
\begin{enumerate}[label=(\Alph*),leftmargin=*,itemsep=0.2em]
\item You could eat the corn all day long.
\item The corn kernels would burn.
\item The kernels will expand from heating.
\item The corn would have more kernels.
\end{enumerate}
\par
\begingroup
\setlength{\parskip}{0pt}
\setlength{\parindent}{0pt}
\textbf{Review}\par
Answer: B\par
Confidence: medium\par
Jawaban diambil dari kunci belajar yang diberikan pengguna dan belum diverifikasi terhadap audio/transkrip resmi.\par
\endgroup
\par
\vspace{0.25\baselineskip}
\noindent\textbf{37.}
\begin{enumerate}[label=(\Alph*),leftmargin=*,itemsep=0.2em]
\item It's many times larger.
\item It's comparatively small.
\item It's hard to measure.
\item It grows steadily.
\end{enumerate}
\par
\begingroup
\setlength{\parskip}{0pt}
\setlength{\parindent}{0pt}
\textbf{Review}\par
Answer: A\par
Confidence: medium\par
Jawaban diambil dari kunci belajar yang diberikan pengguna dan belum diverifikasi terhadap audio/transkrip resmi.\par
\endgroup
\par
\vspace{0.25\baselineskip}
\subsection*{Part C}
DIRECTIONS: In Part C you will hear short lectures and extended conversations. At the end of each, you will be asked several questions. Each lecture or conversation and each question will be spoken only one time. For this reason, you must listen carefully to understand what each speaker says. After you hear a question, read the four choices and select the one that best answers the question the speaker asked. Then, on your answer sheet, find the number of the question and blacken the space that contains the letter for the answer you have chosen. Answer all questions according to what is stated or implied in the lecture or conversation.\par
\begin{exampleblock}
Listen to this sample talk.\par
You will hear:\par
Now listen, to the following example.\par
You will hear:\par
You will read:\par
(A) By cars and carriages\par
(B) By bicycles, trains, and carriages\par
(C) On foot and by boat\par
(D) On board ships and trains\par
The best answer to the question "According to the speaker, how did people travel before the invention of the automobile?" is (B), "By bicycles, trains, and carriages." Therefore, the correct answer is (B). You are not allowed to make notes during the test.\par
\end{exampleblock}
\begin{exampleblock}
\textbf{Transkrip rekonstruksi AI (draft; bukan resmi)}\par
Speaker: Welcome to the course. Before we begin the first lecture, I want to review several course policies. Each week's reading and written work must be completed during the week in which it is assigned. Please do not save assignments until the end of the term. If you are absent for more than a few days because of illness or another serious matter, send a notice to my office or to the department. If you miss a test or assignment deadline, you must inform me of the circumstances as soon as possible. Makeup work may be arranged only when I understand the reason for the absence. Otherwise, late assignments will receive lower grades.\par
\end{exampleblock}
\noindent\textbf{38.}
\begin{enumerate}[label=(\Alph*),leftmargin=*,itemsep=0.2em]
\item Course assignments
\item Course policies
\item Class participation
\item Writing projects
\end{enumerate}
\par
\begingroup
\setlength{\parskip}{0pt}
\setlength{\parindent}{0pt}
\textbf{Review}\par
Answer: B\par
Confidence: medium\par
Jawaban diambil dari kunci belajar yang diberikan pengguna dan belum diverifikasi terhadap audio/transkrip resmi.\par
\endgroup
\par
\vspace{0.25\baselineskip}
\noindent\textbf{39.}
\begin{enumerate}[label=(\Alph*),leftmargin=*,itemsep=0.2em]
\item By the time the paper is due
\item Over a three-day period
\item Before the end of the term
\item During the week assigned
\end{enumerate}
\par
\begingroup
\setlength{\parskip}{0pt}
\setlength{\parindent}{0pt}
\textbf{Review}\par
Answer: D\par
Confidence: medium\par
Jawaban diambil dari kunci belajar yang diberikan pengguna dan belum diverifikasi terhadap audio/transkrip resmi.\par
\endgroup
\par
\vspace{0.25\baselineskip}
\noindent\textbf{40.}
\begin{enumerate}[label=(\Alph*),leftmargin=*,itemsep=0.2em]
\item Assign more homework
\item Take class attendance
\item Give a makeup test
\item Send in a notice
\end{enumerate}
\par
\begingroup
\setlength{\parskip}{0pt}
\setlength{\parindent}{0pt}
\textbf{Review}\par
Answer: D\par
Confidence: medium\par
Jawaban diambil dari kunci belajar yang diberikan pengguna dan belum diverifikasi terhadap audio/transkrip resmi.\par
\endgroup
\par
\vspace{0.25\baselineskip}
\noindent\textbf{41.}
\begin{enumerate}[label=(\Alph*),leftmargin=*,itemsep=0.2em]
\item By informing the speaker of the circumstances
\item By talking to their advisers about the class
\item By handing in assignments
\item By taking makeup tests they have missed
\end{enumerate}
\par
\begingroup
\setlength{\parskip}{0pt}
\setlength{\parindent}{0pt}
\textbf{Review}\par
Answer: A\par
Confidence: medium\par
Jawaban diambil dari kunci belajar yang diberikan pengguna dan belum diverifikasi terhadap audio/transkrip resmi.\par
\endgroup
\par
\vspace{0.25\baselineskip}
\noindent\textbf{42.}
\begin{enumerate}[label=(\Alph*),leftmargin=*,itemsep=0.2em]
\item Special circumstances may arise at any time.
\item No reasonable excuse will ever be accepted.
\item Assignment grades will be lowered otherwise.
\item Urgent matters must be taken care of in any event.
\end{enumerate}
\par
\begingroup
\setlength{\parskip}{0pt}
\setlength{\parindent}{0pt}
\textbf{Review}\par
Answer: C\par
Confidence: medium\par
Jawaban diambil dari kunci belajar yang diberikan pengguna dan belum diverifikasi terhadap audio/transkrip resmi.\par
\endgroup
\par
\vspace{0.25\baselineskip}
\begin{exampleblock}
\textbf{Transkrip rekonstruksi AI (draft; bukan resmi)}\par
Speaker: Today we will discuss how diseases spread from one person to another. A common example is the flu. When a person with the flu coughs, sneezes, or even breathes, tiny particles carrying viruses may be released into the air. Other people can become infected when they breathe in those particles. Bacterial infections can also spread in similar ways, although some diseases may require direct contact or contaminated surfaces. One important point is that people may transmit a virus even before they realize they are ill. For that reason, all infected individuals can potentially spread the disease, not only those who already show clear symptoms.\par
\end{exampleblock}
\noindent\textbf{43.}
\begin{enumerate}[label=(\Alph*),leftmargin=*,itemsep=0.2em]
\item The spread of disease
\item The flu epidemic
\item Flu symptoms
\item Disease variations
\end{enumerate}
\par
\begingroup
\setlength{\parskip}{0pt}
\setlength{\parindent}{0pt}
\textbf{Review}\par
Answer: A\par
Confidence: medium\par
Jawaban diambil dari kunci belajar yang diberikan pengguna dan belum diverifikasi terhadap audio/transkrip resmi.\par
\endgroup
\par
\vspace{0.25\baselineskip}
\noindent\textbf{44.}
\begin{enumerate}[label=(\Alph*),leftmargin=*,itemsep=0.2em]
\item The multiplying of viruses
\item Viruses transported by air
\item Drops of contaminated water
\item The regular breathing of air
\end{enumerate}
\par
\begingroup
\setlength{\parskip}{0pt}
\setlength{\parindent}{0pt}
\textbf{Review}\par
Answer: B\par
Confidence: medium\par
Jawaban diambil dari kunci belajar yang diberikan pengguna dan belum diverifikasi terhadap audio/transkrip resmi.\par
\endgroup
\par
\vspace{0.25\baselineskip}
\noindent\textbf{45.}
\begin{enumerate}[label=(\Alph*),leftmargin=*,itemsep=0.2em]
\item By adhering to droplets in the air
\item Only by tactile contact
\item Similarly to viral infections
\item Frequently through complications
\end{enumerate}
\par
\begingroup
\setlength{\parskip}{0pt}
\setlength{\parindent}{0pt}
\textbf{Review}\par
Answer: C\par
Confidence: medium\par
Jawaban diambil dari kunci belajar yang diberikan pengguna dan belum diverifikasi terhadap audio/transkrip resmi.\par
\endgroup
\par
\vspace{0.25\baselineskip}
\noindent\textbf{46.}
\begin{enumerate}[label=(\Alph*),leftmargin=*,itemsep=0.2em]
\item People who have overt flu symptoms
\item Humans who seek contact
\item Persons who are not self-aware
\item All those infected with a virus
\end{enumerate}
\par
\begingroup
\setlength{\parskip}{0pt}
\setlength{\parindent}{0pt}
\textbf{Review}\par
Answer: D\par
Confidence: medium\par
Jawaban diambil dari kunci belajar yang diberikan pengguna dan belum diverifikasi terhadap audio/transkrip resmi.\par
\endgroup
\par
\vspace{0.25\baselineskip}
\begin{exampleblock}
\textbf{Transkrip rekonstruksi AI (draft; bukan resmi)}\par
Speaker: Opinion researchers often study why people choose certain goods and services. In recent surveys, many consumers have said that time is one of their most valuable resources. This helps explain the popularity of fast-food restaurants, household services, and disposable goods, since these products can save time. Researchers also find that appearance, or how people look, strongly affects many buying decisions, especially in areas such as clothing, cosmetics, and exercise programs. But not every consumer choice is practical or appearance-related. Many people spend money on hobbies such as gardening, decorating, cooking, or crafts because these activities allow them to express personal creativity.\par
\end{exampleblock}
\noindent\textbf{47.}
\begin{enumerate}[label=(\Alph*),leftmargin=*,itemsep=0.2em]
\item Product manufacturers
\item Opinion researchers
\item Restaurant owners
\item Commodity traders
\end{enumerate}
\par
\begingroup
\setlength{\parskip}{0pt}
\setlength{\parindent}{0pt}
\textbf{Review}\par
Answer: B\par
Confidence: medium\par
Jawaban diambil dari kunci belajar yang diberikan pengguna dan belum diverifikasi terhadap audio/transkrip resmi.\par
\endgroup
\par
\vspace{0.25\baselineskip}
\noindent\textbf{48.}
\begin{enumerate}[label=(\Alph*),leftmargin=*,itemsep=0.2em]
\item Time
\item Money
\item Fast-food restaurants
\item Household services
\end{enumerate}
\par
\begingroup
\setlength{\parskip}{0pt}
\setlength{\parindent}{0pt}
\textbf{Review}\par
Answer: A\par
Confidence: medium\par
Jawaban diambil dari kunci belajar yang diberikan pengguna dan belum diverifikasi terhadap audio/transkrip resmi.\par
\endgroup
\par
\vspace{0.25\baselineskip}
\noindent\textbf{49.}
\begin{enumerate}[label=(\Alph*),leftmargin=*,itemsep=0.2em]
\item Lawns
\item Fast-drying paints
\item Looks
\item Disposable goods
\end{enumerate}
\par
\begingroup
\setlength{\parskip}{0pt}
\setlength{\parindent}{0pt}
\textbf{Review}\par
Answer: C\par
Confidence: medium\par
Jawaban diambil dari kunci belajar yang diberikan pengguna dan belum diverifikasi terhadap audio/transkrip resmi.\par
\endgroup
\par
\vspace{0.25\baselineskip}
\noindent\textbf{50.}
\begin{enumerate}[label=(\Alph*),leftmargin=*,itemsep=0.2em]
\item To stay healthy
\item To improve their appearance
\item To become energetic
\item To satisfy personal creativity
\end{enumerate}
\par
\begingroup
\setlength{\parskip}{0pt}
\setlength{\parindent}{0pt}
\textbf{Review}\par
Answer: D\par
Confidence: medium\par
Jawaban diambil dari kunci belajar yang diberikan pengguna dan belum diverifikasi terhadap audio/transkrip resmi.\par
\endgroup
\par
\vspace{0.25\baselineskip}
\section*{Structure and Written Expression}
\subsection*{Sentence Completion}
This section is designed to test your ability to recognize language structures that are appropriate in standard written English. The questions in this section belong to two types, each of which has special directions.\par
DIRECTIONS: Questions 1-15 are partial sentences. Below each sentence you will see four words or phrases, marked (A), (B), (C), and (D). Select the one word or phrase that best completes the sentence. Then, on your answer sheet, find the number of the question and fill in the space that contains the letter for the answer you have chosen. Fill in the space completely.\par
\begin{exampleblock}
\textbf{EXAMPLE I}\par
\noindent Drying flowers is the best way......them.
\begin{enumerate}[label=(\Alph*),leftmargin=*,itemsep=0.2em]
\item to preserve
\item by preserving
\item preserve
\item preserved
\end{enumerate}
\vspace{0.25\baselineskip}
The sentence should state, "Drying flowers is the best way to preserve them." Therefore, the correct answer is (A).\par
\end{exampleblock}
\begin{exampleblock}
\textbf{EXAMPLE II}\par
\noindent Many American universities......as small, private colleges.
\begin{enumerate}[label=(\Alph*),leftmargin=*,itemsep=0.2em]
\item begun
\item beginning
\item began
\item for the beginning
\end{enumerate}
\vspace{0.25\baselineskip}
The sentence should state, "Many American universities began as small, private colleges." Therefore, the correct answer is (C). After you read the directions, begin work on the questions.\par
\end{exampleblock}
\noindent\textbf{1.} Tennessee has about 140 newspapers,......25 are issued daily.
\begin{enumerate}[label=(\Alph*),leftmargin=*,itemsep=0.2em]
\item about which
\item of which about
\item which are about
\item which is about
\end{enumerate}
\par
\begingroup
\setlength{\parskip}{0pt}
\setlength{\parindent}{0pt}
\textbf{Review}\par
Answer: B\par
Confidence: high\par
Grammar rule: Quantifier in a nonrestrictive relative clause with of which.\par
Setelah koma, kalimat membutuhkan relative clause nonrestrictive untuk menyatakan jumlah bagian dari `140 newspapers`; pola yang benar adalah `of which about 25 are issued daily`.\par
\endgroup
\par
\vspace{0.25\baselineskip}
\noindent\textbf{2.} When consumers cannot have everything they want, they have to choose......most.
\begin{enumerate}[label=(\Alph*),leftmargin=*,itemsep=0.2em]
\item they want what
\item what they want
\item they want it
\item that they want
\end{enumerate}
\par
\begingroup
\setlength{\parskip}{0pt}
\setlength{\parindent}{0pt}
\textbf{Review}\par
Answer: B\par
Confidence: high\par
Grammar rule: Embedded noun clause after choose.\par
Setelah `choose`, dibutuhkan object clause sebagai sesuatu yang dipilih; `what they want` membentuk klausa benda yang lengkap dan alami sebelum `most`.\par
\endgroup
\par
\vspace{0.25\baselineskip}
\noindent\textbf{3.} The temperature of an object rises when......into it.
\begin{enumerate}[label=(\Alph*),leftmargin=*,itemsep=0.2em]
\item heat flow
\item flows hot
\item heat flows
\item hot flow
\end{enumerate}
\par
\begingroup
\setlength{\parskip}{0pt}
\setlength{\parindent}{0pt}
\textbf{Review}\par
Answer: C\par
Confidence: high\par
Grammar rule: Simple present clause after when.\par
Sesudah `when`, diperlukan klausa lengkap dengan subject + verb; `heat flows` tepat karena `heat` menjadi subjek dan `flows` menjadi verb tunggal.\par
\endgroup
\par
\vspace{0.25\baselineskip}
\noindent\textbf{4.} From ancient times, people......their land, air and water.
\begin{enumerate}[label=(\Alph*),leftmargin=*,itemsep=0.2em]
\item always have polluting
\item always have pollution
\item have always polluted
\item pollution always has
\end{enumerate}
\par
\begingroup
\setlength{\parskip}{0pt}
\setlength{\parindent}{0pt}
\textbf{Review}\par
Answer: C\par
Confidence: high\par
Grammar rule: Present perfect with adverb placement.\par
Keterangan `From ancient times` cocok dengan present perfect untuk aksi sejak masa lalu sampai kini; posisi adverb yang benar adalah `have always polluted`.\par
\endgroup
\par
\vspace{0.25\baselineskip}
\noindent\textbf{5.} Jean Fragonard, was a French artist......portraits of children.
\begin{enumerate}[label=(\Alph*),leftmargin=*,itemsep=0.2em]
\item whose paintings
\item who has painted
\item who painted
\item whose painted
\end{enumerate}
\par
\begingroup
\setlength{\parskip}{0pt}
\setlength{\parindent}{0pt}
\textbf{Review}\par
Answer: C\par
Confidence: high\par
Grammar rule: Relative clause with who + past verb.\par
Setelah `a French artist`, dibutuhkan relative clause yang menjelaskan orang tersebut; `who painted portraits of children` paling tepat dan lengkap secara struktur.\par
\endgroup
\par
\vspace{0.25\baselineskip}
\noindent\textbf{6.} Overharvesting brought North American alligators, to......in their natural
\begin{enumerate}[label=(\Alph*),leftmargin=*,itemsep=0.2em]
\item nearly extinct
\item near extinction
\item extinct near
\item extinction nearly
\end{enumerate}
\par
\begingroup
\setlength{\parskip}{0pt}
\setlength{\parindent}{0pt}
\textbf{Review}\par
Answer: B\par
Confidence: medium\par
Grammar rule: Fixed expression with bring to + noun.\par
Pola yang umum adalah `bring ... to near extinction`, yaitu `to` diikuti noun phrase, bukan adjective. Sumber kalimat tampak terpotong, tetapi secara struktur pilihan ini paling sesuai.\par
Needs manual review: true\par
\endgroup
\par
\vspace{0.25\baselineskip}
\noindent\textbf{7.} The Ford Foundation was established in 1936 to advance human well-being by......funds for education.
\begin{enumerate}[label=(\Alph*),leftmargin=*,itemsep=0.2em]
\item contribute
\item contribution
\item to contribute
\item contributing
\end{enumerate}
\par
\begingroup
\setlength{\parskip}{0pt}
\setlength{\parindent}{0pt}
\textbf{Review}\par
Answer: D\par
Confidence: high\par
Grammar rule: By + gerund to express means.\par
Sesudah preposisi `by`, verb harus berbentuk gerund; karena itu `by contributing funds for education` adalah pola yang benar.\par
\endgroup
\par
\vspace{0.25\baselineskip}
\noindent\textbf{8.} Fireweed received its name because it......after a forest fire.
\begin{enumerate}[label=(\Alph*),leftmargin=*,itemsep=0.2em]
\item quick growth
\item grows quickly
\item quickly grown
\item growing quickly
\end{enumerate}
\par
\begingroup
\setlength{\parskip}{0pt}
\setlength{\parindent}{0pt}
\textbf{Review}\par
Answer: B\par
Confidence: high\par
Grammar rule: Simple present verb in a reason clause.\par
Setelah `because it`, diperlukan predicate yang lengkap; `grows quickly` memberi verb yang benar dan adverb `quickly` memodifikasi cara tumbuhnya.\par
\endgroup
\par
\vspace{0.25\baselineskip}
\noindent\textbf{9.} Florida's long coastline and warm weather......swimmers to its sandy shores.
\begin{enumerate}[label=(\Alph*),leftmargin=*,itemsep=0.2em]
\item attracts
\item attract
\item they attract
\item is attracted by
\end{enumerate}
\par
\begingroup
\setlength{\parskip}{0pt}
\setlength{\parindent}{0pt}
\textbf{Review}\par
Answer: B\par
Confidence: high\par
Grammar rule: Subject-verb agreement with a compound subject.\par
Subjek `Florida's long coastline and warm weather` bersifat jamak karena dua unsur digabung dengan `and`, jadi verb yang tepat adalah `attract`.\par
\endgroup
\par
\vspace{0.25\baselineskip}
\noindent\textbf{10.} Amazon pygmies consider their songs......part of their culture.
\begin{enumerate}[label=(\Alph*),leftmargin=*,itemsep=0.2em]
\item an important and extremely
\item as extreme and important
\item an extremely important
\item as extreme importance
\end{enumerate}
\par
\begingroup
\setlength{\parskip}{0pt}
\setlength{\parindent}{0pt}
\textbf{Review}\par
Answer: C\par
Confidence: high\par
Grammar rule: Object complement noun phrase after consider.\par
Pola `consider + object + noun phrase` membutuhkan pelengkap seperti `an extremely important part`; adjective `extremely important` menerangkan noun `part` dengan urutan yang benar.\par
\endgroup
\par
\vspace{0.25\baselineskip}
\noindent\textbf{11.} Psychologists define anxiety as a feeling of dread, apprehension, or......
\begin{enumerate}[label=(\Alph*),leftmargin=*,itemsep=0.2em]
\item afraid
\item be afraid
\item having fear
\item fear
\end{enumerate}
\par
\begingroup
\setlength{\parskip}{0pt}
\setlength{\parindent}{0pt}
\textbf{Review}\par
Answer: D\par
Confidence: high\par
Grammar rule: Parallel nouns after or.\par
Setelah `a feeling of dread, apprehension, or ...`, unsur-unsurnya harus paralel sebagai noun; `fear` sejajar dengan `dread` dan `apprehension`.\par
\endgroup
\par
\vspace{0.25\baselineskip}
\noindent\textbf{12.} Ancient nations have used......on their emblems and flags for thousands of years.
\begin{enumerate}[label=(\Alph*),leftmargin=*,itemsep=0.2em]
\item the same symbols as
\item the same symbols
\item symbols the same as
\item symbols as the same
\end{enumerate}
\par
\begingroup
\setlength{\parskip}{0pt}
\setlength{\parindent}{0pt}
\textbf{Review}\par
Answer: B\par
Confidence: high\par
Grammar rule: Direct object noun phrase after used.\par
Verb `have used` membutuhkan object langsung; `the same symbols` sudah lengkap, sedangkan pilihan lain menambah `as` tanpa pasangan pembanding yang diperlukan.\par
\endgroup
\par
\vspace{0.25\baselineskip}
\noindent\textbf{13.} For years experts......the effect of coaching and preparatory courses on test scores.
\begin{enumerate}[label=(\Alph*),leftmargin=*,itemsep=0.2em]
\item are examining
\item had been examined
\item have been examining.
\item having been examined
\end{enumerate}
\par
\begingroup
\setlength{\parskip}{0pt}
\setlength{\parindent}{0pt}
\textbf{Review}\par
Answer: C\par
Confidence: high\par
Grammar rule: Present perfect continuous with for + period of time.\par
Keterangan `For years` menandakan aktivitas yang berlangsung selama periode waktu sampai sekarang, sehingga `have been examining` adalah tense yang paling tepat.\par
\endgroup
\par
\vspace{0.25\baselineskip}
\noindent\textbf{14.} The nitrogen cycle......of nitrogen through the atmosphere, hydrosphere, and lithosphere.
\begin{enumerate}[label=(\Alph*),leftmargin=*,itemsep=0.2em]
\item the circulation is
\item is the circulation
\item it is the circulation
\item is it the circulation
\end{enumerate}
\par
\begingroup
\setlength{\parskip}{0pt}
\setlength{\parindent}{0pt}
\textbf{Review}\par
Answer: B\par
Confidence: high\par
Grammar rule: Subject-complement sentence with be.\par
Kalimat definisi membutuhkan pola `subject + be + complement`; jadi `The nitrogen cycle is the circulation of nitrogen ...` adalah susunan yang benar.\par
\endgroup
\par
\vspace{0.25\baselineskip}
\noindent\textbf{15.} While working as a clerk, Edison spent much of his time......the stock ticker.
\begin{enumerate}[label=(\Alph*),leftmargin=*,itemsep=0.2em]
\item he studied
\item to study
\item on study
\item studying
\end{enumerate}
\par
\begingroup
\setlength{\parskip}{0pt}
\setlength{\parindent}{0pt}
\textbf{Review}\par
Answer: D\par
Confidence: high\par
Grammar rule: Spend time + gerund.\par
Setelah pola `spent much of his time`, bentuk yang benar adalah gerund; karena itu `studying the stock ticker` paling sesuai.\par
\endgroup
\par
\vspace{0.25\baselineskip}
\subsection*{Error Identification}
DIRECTIONS: In questions 16-40 every sentence has four words or phrases that are underlined. The four underlined portions of each sentence are marked (A), (B), (C), and (D). Identify the one word or phrase that makes the sentence incorrect. Then, on your answer sheet, find the number of the question and fill in the space that contains the letter for the answer you have chosen. Fill in the space completely.\par
\begin{exampleblock}
\textbf{EXAMPLE I}\par
Christopher Columbus \uline{has sailed}(A) from Europe in 1492 and \uline{discovered a}(B) \uline{new land}(C) he \uline{thought to}(D) be India.\par
The sentence should state, "Christopher Columbus sailed from Europe in 1492 and discovered a new land he thought to be India." Therefore, you should choose answer (A).\par
\end{exampleblock}
\begin{exampleblock}
\textbf{EXAMPLE II}\par
\uline{As the roles}(A) of people in society change, \uline{so does}(B) the \uline{rules of}(C) conduct in certain \uline{situations}(D).\par
The sentence should state, "As the roles of people in society change, so do the rules of\par
conduct in certain situations." Therefore, you should choose answer (B).\par
\end{exampleblock}
\noindent\textbf{16.} James Maxwell \uline{based}(A) his \uline{work on}(B) the \uline{discoveries}(C) of the English \uline{physical}(D) Michael Faraday.
\par
\begingroup
\setlength{\parskip}{0pt}
\setlength{\parindent}{0pt}
\textbf{Review}\par
Answer: D\par
Confidence: high\par
Grammar rule: Noun form after adjective.\par
Correct form: physicist\par
Bagian (D) salah karena setelah kata sifat `English` dibutuhkan nomina untuk menyebut profesi. Bentuk yang benar adalah `physicist`.\par
\endgroup
\par
\vspace{\questionblockgap}
\noindent\textbf{17.} The behavior of animals \uline{appears to}(A) depend on \uline{patterns of}(B) \uline{reactions which}(C) \uline{they are}(D) born.
\par
\begingroup
\setlength{\parskip}{0pt}
\setlength{\parindent}{0pt}
\textbf{Review}\par
Answer: D\par
Confidence: high\par
Grammar rule: Required preposition in relative clause.\par
Correct form: they are born with\par
Bagian (D) salah karena frasa `born` dalam makna ini memerlukan preposisi `with`. Bentuk yang benar adalah `they are born with`.\par
\endgroup
\par
\vspace{\questionblockgap}
\noindent\textbf{18.} \uline{For centuries}(A), people \uline{have wondered}(B) why \uline{have they}(C) particular dreams while \uline{they sleep}(D).
\par
\begingroup
\setlength{\parskip}{0pt}
\setlength{\parindent}{0pt}
\textbf{Review}\par
Answer: C\par
Confidence: high\par
Grammar rule: Indirect question word order.\par
Correct form: they have\par
Bagian (C) salah karena indirect question tidak memakai inversi seperti pertanyaan langsung. Bentuk yang benar adalah `they have`.\par
\endgroup
\par
\vspace{\questionblockgap}
\noindent\textbf{19.} \uline{After the}(A) Roman Empire \uline{has collapsed}(B), Europe had \uline{no regular}(C) \uline{postal}(D) service.
\par
\begingroup
\setlength{\parskip}{0pt}
\setlength{\parindent}{0pt}
\textbf{Review}\par
Answer: B\par
Confidence: high\par
Grammar rule: Verb tense consistency in past narration.\par
Correct form: collapsed\par
Bagian (B) salah karena kalimat naratif lampau ini tidak tepat memakai present perfect. Bentuk yang benar adalah `collapsed`.\par
\endgroup
\par
\vspace{\questionblockgap}
\noindent\textbf{20.} Most frequently, asphalt \uline{is used}(A) \uline{to pave}(B) streets, \uline{avenues}(C), highways, and \uline{airport}(D).
\par
\begingroup
\setlength{\parskip}{0pt}
\setlength{\parindent}{0pt}
\textbf{Review}\par
Answer: D\par
Confidence: high\par
Grammar rule: Parallel plural nouns in a series.\par
Correct form: airports\par
Bagian (D) salah karena unsur-unsur dalam deret harus sejajar dalam bentuk jamak: streets, avenues, highways, dan airports. Bentuk yang benar adalah `airports`.\par
\endgroup
\par
\vspace{\questionblockgap}
\noindent\textbf{21.} Samuel Coleridge was \uline{poet}(A) and philosopher of \uline{the English}(B) \uline{romantic}(C) movement in \uline{the early}(D) 1800s.
\par
\begingroup
\setlength{\parskip}{0pt}
\setlength{\parindent}{0pt}
\textbf{Review}\par
Answer: A\par
Confidence: high\par
Grammar rule: Missing article before singular count noun.\par
Correct form: a poet\par
Bagian (A) salah karena nomina tunggal yang dapat dihitung memerlukan artikel. Bentuk yang benar adalah `a poet`.\par
\endgroup
\par
\vspace{\questionblockgap}
\noindent\textbf{22.} When one person \uline{intentional}(A) takes the property of \uline{another}(B) without legal \uline{justification}(C), the crime is \uline{called}(D) theft.
\par
\begingroup
\setlength{\parskip}{0pt}
\setlength{\parindent}{0pt}
\textbf{Review}\par
Answer: A\par
Confidence: high\par
Grammar rule: Adverb form modifying a verb.\par
Correct form: intentionally\par
Bagian (A) salah karena kata yang menerangkan verba `takes` harus berbentuk adverb. Bentuk yang benar adalah `intentionally`.\par
\endgroup
\par
\vspace{\questionblockgap}
\noindent\textbf{23.} \uline{Computers}(A) programs, \uline{catalogues}(B), directories, and \uline{collections}(C) of data are \uline{protected}(D) by copyright laws.
\par
\begingroup
\setlength{\parskip}{0pt}
\setlength{\parindent}{0pt}
\textbf{Review}\par
Answer: A\par
Confidence: high\par
Grammar rule: Noun modifier form.\par
Correct form: Computer\par
Bagian (A) salah karena nomina yang berfungsi sebagai modifier biasanya berbentuk tunggal. Bentuk yang benar adalah `Computer` dalam frasa `Computer programs`.\par
\endgroup
\par
\vspace{\questionblockgap}
\noindent\textbf{24.} The water temperature in \uline{a spring}(A) depends \uline{on that}(B) of the soil through \uline{where}(C) the water \uline{flows}(D).
\par
\begingroup
\setlength{\parskip}{0pt}
\setlength{\parindent}{0pt}
\textbf{Review}\par
Answer: C\par
Confidence: high\par
Grammar rule: Relative pronoun after preposition.\par
Correct form: which\par
Bagian (C) salah karena setelah preposisi `through` harus digunakan relative pronoun `which`, bukan `where`. Bentuk yang benar adalah `which`.\par
\endgroup
\par
\vspace{\questionblockgap}
\noindent\textbf{25.} Tree squirrels are active, \uline{noisy}(A), and \uline{lively}(B) \uline{animals that}(C) make \uline{its home}(D) in tree trunks.
\par
\begingroup
\setlength{\parskip}{0pt}
\setlength{\parindent}{0pt}
\textbf{Review}\par
Answer: D\par
Confidence: high\par
Grammar rule: Pronoun agreement.\par
Correct form: their homes\par
Bagian (D) salah karena pronoun harus sesuai dengan subjek jamak `Tree squirrels`. Bentuk yang benar adalah `their homes`.\par
\endgroup
\par
\vspace{\questionblockgap}
\noindent\textbf{26.} Botanists \uline{have determined}(A) that \uline{there is}(B) more than 60 \uline{species}(C) of \uline{sunflowers}(D).
\par
\begingroup
\setlength{\parskip}{0pt}
\setlength{\parindent}{0pt}
\textbf{Review}\par
Answer: B\par
Confidence: high\par
Grammar rule: Subject-verb agreement with existential there.\par
Correct form: there are\par
Bagian (B) salah karena setelah existential `there` dengan subjek jamak `more than 60 species` harus dipakai verba jamak. Bentuk yang benar adalah `there are`.\par
\endgroup
\par
\vspace{\questionblockgap}
\noindent\textbf{27.} Pure sodium \uline{immediately}(A) \uline{combines}(B) with oxygen when \uline{is exposed}(C) \uline{to air}(D).
\par
\begingroup
\setlength{\parskip}{0pt}
\setlength{\parindent}{0pt}
\textbf{Review}\par
Answer: C\par
Confidence: high\par
Grammar rule: Missing subject in adverb clause.\par
Correct form: it is exposed\par
Bagian (C) salah karena anak kalimat setelah `when` memerlukan subjek. Bentuk yang benar adalah `it is exposed`.\par
\endgroup
\par
\vspace{\questionblockgap}
\noindent\textbf{28.} Special education is intended \uline{help}(A) \uline{both}(B) handicapped and gifted children to \uline{reach}(C) their \uline{learning}(D) potentials.
\par
\begingroup
\setlength{\parskip}{0pt}
\setlength{\parindent}{0pt}
\textbf{Review}\par
Answer: A\par
Confidence: high\par
Grammar rule: Infinitive after intended.\par
Correct form: to help\par
Bagian (A) salah karena pola yang benar adalah `be intended to + verb`. Bentuk yang benar adalah `to help`.\par
\endgroup
\par
\vspace{\questionblockgap}
\noindent\textbf{29.} \uline{Trading}(A) fairs \uline{held}(B) in Antwerp during the 1300s \uline{brought}(C) \uline{famous}(D) to the city.
\par
\begingroup
\setlength{\parskip}{0pt}
\setlength{\parindent}{0pt}
\textbf{Review}\par
Answer: D\par
Confidence: high\par
Grammar rule: Noun vs adjective form.\par
Correct form: fame\par
Bagian (D) salah karena setelah verba `brought` dibutuhkan nomina sebagai objek, bukan adjective. Bentuk yang benar adalah `fame`.\par
\endgroup
\par
\vspace{\questionblockgap}
\noindent\textbf{30.} By 1850, California had \uline{a sufficiently}(A) \uline{large population}(B) \uline{to recognized}(C) \uline{as a state}(D).
\par
\begingroup
\setlength{\parskip}{0pt}
\setlength{\parindent}{0pt}
\textbf{Review}\par
Answer: C\par
Confidence: high\par
Grammar rule: Infinitive and passive form.\par
Correct form: to be recognized\par
Bagian (C) salah karena setelah `to` harus dipakai bentuk infinitive yang tepat, dan maknanya menuntut bentuk pasif. Bentuk yang benar adalah `to be recognized`.\par
\endgroup
\par
\vspace{\questionblockgap}
\noindent\textbf{31.} Gingham is \uline{a fabric}(A) used to make \uline{dresses}(B), \uline{curtains}(C), and \uline{furnitures}(D) covers.
\par
\begingroup
\setlength{\parskip}{0pt}
\setlength{\parindent}{0pt}
\textbf{Review}\par
Answer: D\par
Confidence: high\par
Grammar rule: Uncountable noun form.\par
Correct form: furniture\par
Bagian (D) salah karena `furniture` adalah uncountable noun dan tidak memakai bentuk jamak berakhiran `-s`. Bentuk yang benar adalah `furniture`.\par
\endgroup
\par
\vspace{\questionblockgap}
\noindent\textbf{32.} Plato's most \uline{last}(A) contribution to mathematics was his \uline{insistence}(B) on \uline{using}(C) \uline{reasoning}(D) in geometry.
\par
\begingroup
\setlength{\parskip}{0pt}
\setlength{\parindent}{0pt}
\textbf{Review}\par
Answer: A\par
Confidence: high\par
Grammar rule: Adjective choice and form.\par
Correct form: lasting\par
Bagian (A) salah karena kolokasi yang tepat adalah `most lasting contribution`, bukan `most last contribution`. Bentuk yang benar adalah `lasting`.\par
\endgroup
\par
\vspace{\questionblockgap}
\noindent\textbf{33.} Sedatives are \uline{a group}(A) of drugs \uline{that legally}(B) prescribed and should \uline{be taken}(C) \uline{as directed}(D).
\par
\begingroup
\setlength{\parskip}{0pt}
\setlength{\parindent}{0pt}
\textbf{Review}\par
Answer: B\par
Confidence: high\par
Grammar rule: Passive verb in relative clause.\par
Correct form: that are legally\par
Bagian (B) salah karena relative clause pasif memerlukan auxiliary `are` sebelum adverb dan past participle. Bentuk yang benar adalah `that are legally`.\par
\endgroup
\par
\vspace{\questionblockgap}
\noindent\textbf{34.} Maps \uline{that show}(A) \uline{detail}(B) landforms \uline{are commonly}(C) \uline{used}(D) in physical geography and geology.
\par
\begingroup
\setlength{\parskip}{0pt}
\setlength{\parindent}{0pt}
\textbf{Review}\par
Answer: B\par
Confidence: high\par
Grammar rule: Adjective form before noun.\par
Correct form: detailed\par
Bagian (B) salah karena kata yang menerangkan nomina `landforms` harus berbentuk adjective. Bentuk yang benar adalah `detailed`.\par
\endgroup
\par
\vspace{\questionblockgap}
\noindent\textbf{35.} When the minerals \uline{needed}(A) for corn \uline{to grow}(B) \uline{are lack}(C), the husks may \uline{be stunted}(D).
\par
\begingroup
\setlength{\parskip}{0pt}
\setlength{\parindent}{0pt}
\textbf{Review}\par
Answer: C\par
Confidence: high\par
Grammar rule: Adjective complement after be.\par
Correct form: are lacking\par
Bagian (C) salah karena setelah verba `are` dibutuhkan adjective atau participle, bukan nomina atau verba dasar `lack`. Bentuk yang benar adalah `are lacking`.\par
\endgroup
\par
\vspace{\questionblockgap}
\noindent\textbf{36.} Rembrandt was \uline{forced}(A) to \uline{declare}(B) bankruptcy in 1656, and his \uline{possessions}(C) were \uline{sale}(D).
\par
\begingroup
\setlength{\parskip}{0pt}
\setlength{\parindent}{0pt}
\textbf{Review}\par
Answer: D\par
Confidence: high\par
Grammar rule: Past participle in passive construction.\par
Correct form: sold\par
Bagian (D) salah karena setelah `were` pada bentuk pasif harus dipakai past participle. Bentuk yang benar adalah `sold`.\par
\endgroup
\par
\vspace{\questionblockgap}
\noindent\textbf{37.} Pearls and \uline{similar}(A) substances \uline{may be}(B) \uline{classified}(C) by \uline{how are}(D) cultivated.
\par
\begingroup
\setlength{\parskip}{0pt}
\setlength{\parindent}{0pt}
\textbf{Review}\par
Answer: D\par
Confidence: high\par
Grammar rule: Embedded clause word order.\par
Correct form: how they are\par
Bagian (D) salah karena klausa setelah preposisi memakai urutan pernyataan, bukan inversi pertanyaan. Bentuk yang benar adalah `how they are`.\par
\endgroup
\par
\vspace{\questionblockgap}
\noindent\textbf{38.} During \uline{winter}(A), grizzly bears \uline{live in}(B) dens, caves or \uline{the other}(C) natural \uline{shelters}(D).
\par
\begingroup
\setlength{\parskip}{0pt}
\setlength{\parindent}{0pt}
\textbf{Review}\par
Answer: C\par
Confidence: high\par
Grammar rule: Article use with plural noun phrase.\par
Correct form: other\par
Bagian (C) salah karena frasa umum jamak tidak memerlukan artikel definit `the` di depan `other natural shelters`. Bentuk yang benar adalah `other`.\par
\endgroup
\par
\vspace{\questionblockgap}
\noindent\textbf{39.} The neck of a \uline{classical}(A) guitar is \uline{wider}(B) than \uline{those}(C) of a steel-\uline{string}(D) guitar.
\par
\begingroup
\setlength{\parskip}{0pt}
\setlength{\parindent}{0pt}
\textbf{Review}\par
Answer: C\par
Confidence: high\par
Grammar rule: Pronoun agreement in comparison.\par
Correct form: that\par
Bagian (C) salah karena pembandingnya adalah nomina tunggal `neck`, sehingga harus memakai pronoun tunggal. Bentuk yang benar adalah `that`.\par
\endgroup
\par
\vspace{\questionblockgap}
\noindent\textbf{40.} \uline{Insufficient}(A) oxygen causes lactic acid to \uline{built up}(B) in the \uline{muscles}(C) of long-\uline{distance}(D) runners.
\par
\begingroup
\setlength{\parskip}{0pt}
\setlength{\parindent}{0pt}
\textbf{Review}\par
Answer: B\par
Confidence: high\par
Grammar rule: Base verb after infinitive.\par
Correct form: build up\par
Bagian (B) salah karena setelah `to` harus dipakai verba bentuk dasar, bukan past participle. Bentuk yang benar adalah `build up`.\par
\endgroup
\par
\vspace{\questionblockgap}
\section*{Reading Comprehension}
\begin{center}
\textbf{Time 55 minutes}
\end{center}
DIRECTIONS: In this section you will read several passages. Each passage is followed by a series of questions. For questions 1-50, you need to select the best answer, (A), (B), (C), or (D), to each question. Then, on your answer sheet, find the number of the question and fill in the space that contains the letter of the answer you have selected. Fill in the space completely. Answer all questions following a passage on the basis of what is stated or implied in the passage.\par
\begin{exampleblock}
\small
\sloppy
Read the following passage:\par
A tomahawk is a small ax used as a tool and a weapon by the North American Indian\par
tribes. An average tomahawk was not very long and did not weigh a great deal. Originally,\par
the head of the tomahawk was made of a shaped stone or an animal bone and was mounted on\par
Line a wooden handle. After the arrival of the European settlers, the Indians began to use toma-\par
hawks with iron heads. Indian males and females of all ages used tomahawks to chop and cut wood, pound stakes into the ground to put up wigwams, and do many other chores. Indian warriors relied on tomahawks as weapons and even threw them at their enemies. Some types of tomahawks were used in religious ceremonies. Contemporary American idioms reflect\par
this aspect of American heritage.\par
\end{exampleblock}
\begin{exampleblock}
\textbf{EXAMPLE I}\par
\noindent Early tomahawk heads were made of
\begin{enumerate}[label=(\Alph*),leftmargin=*,itemsep=0.2em]
\item stone or bone
\item wood or sticks
\item European iron
\item religious weapons
\end{enumerate}
\vspace{0.25\baselineskip}
According to the passage, early tomahawk heads were made of stone or bone. Therefore, the correct answer is (A).\par
\end{exampleblock}
\begin{exampleblock}
\textbf{EXAMPLE II}\par
\noindent How has the Indian use of tomahawks affected American daily life today?
\begin{enumerate}[label=(\Alph*),leftmargin=*,itemsep=0.2em]
\item Tomahawks are still used as weapons.
\item Tomahawks are used as tools for.certain jobs.
\item Contemporary language refers to tomahawks.
\item Indian tribes cherish tomahawks as heirlooms.
\end{enumerate}
\vspace{0.25\baselineskip}
The passage states, "Contemporary American idioms reflect this aspect of American heritage." The correct answer is (C). After you read the directions, begin work on the questions.\par
\end{exampleblock}
\subsection*{Questions 1-10}
\begingroup
\small
\sloppy
\setlength{\parskip}{0pt}
\setlength{\parindent}{0pt}
During the Middle Ages, societies were based on military relationships, as landowners\par
formed their own foot armies into which they drafted their tenants and hired hands. The in-\par
fantry that fought its way forward against the opposition engaged in heavy ground battles\par
that proved costly in the ratio of losses to wins. These soldiers carried darts, javelins, and\par
slings to be used before closing ranks with the enemy, although their swords and halberds\par
 delivered crushing blows on contact. Such armed forces were active for limited periods of\par
time and had a predominantly defensive function, displayed in hand-to-hand combat.\par
Because this sporadic and untrained organization was ineffective, the ruling classes be-\par
gan to hire mercenaries who were generously compensated for their tasks and subject to con-\par
tractual terms of agreement. The greatest idiosyncrasy of a hired military force was that the\par
troops sometimes deserted their employers if they could bank on a higher remuneration\par
from the opposition. The Swiss pikemen became the best-known mercenaries of the late\par
Middle Ages. In the 1300s, they practically invented a crude body armor of leather and\par
quilted layered head gear with nose and skull plates, ornamented with crests. Their tower\par
shields proved indispensable against a shower of arrows, and their helmets progressed from\par
cone cups to visors hinged at the temples. As their notoriety increased, so did their wages,\par
and eventually they were rounded into military companies that later grew into the basic\par
units in almost all armies. During the same period, the first full-size army of professional\par
soldiers emerged in the Ottoman Empire. What set these troops apart from other contempo-\par
rary armies was that these soldiers remained on duty in peacetime.\par
Companies of mercenaries were employed on a permanent basis in 1445, when King\par
Charles VII created a regular military organization, complete with a designated hierarchy.\par
Gunpowder accelerated the emergence of military tactics and strategy that ultimately af-\par
fected the conceptualization of war on a broad scale. Cannons further widened the gap be-\par
tween the attacking and the defending lineups, and undermined the exclusivity of contact\par
battles.\par
\endgroup
\medskip
\noindent\textbf{1.} What is the main purpose of the passage?
\begin{enumerate}[label=(\Alph*),leftmargin=*,itemsep=0.2em]
\item To distinguish between laborers and mercenaries
\item To change the existing view of the military
\item To cite examples of armor in the Middle Ages
\item To trace the origins of military organization
\end{enumerate}
\par
\begingroup
\setlength{\parskip}{0pt}
\setlength{\parindent}{0pt}
\textbf{Review}\par
Answer: D\par
Confidence: high\par
Evidence refs: lines 1-2, lines 7-10, lines 14-20\par
Evidence quote: "During the Middle Ages, societies were based on military relationships... ; the ruling classes began to hire mercenaries... ; the first full-size army of professional soldiers emerged... Gunpowder accelerated the emergence of military tactics..."\par
Jawaban (D) karena passage menelusuri perkembangan organisasi militer dari pasukan feodal, lalu mercenary, hingga tentara profesional dan perubahan taktik perang.\par
\endgroup
\par
\vspace{0.25\baselineskip}
\noindent\textbf{2.} In line 4, the word "ratio" is closest in meaning to
\begin{enumerate}[label=(\Alph*),leftmargin=*,itemsep=0.2em]
\item quota
\item reason
\item proportion
\item pace
\end{enumerate}
\par
\begingroup
\setlength{\parskip}{0pt}
\setlength{\parindent}{0pt}
\textbf{Review}\par
Answer: C\par
Confidence: high\par
Evidence refs: line 4\par
Evidence quote: "that proved costly in the ratio of losses to wins"\par
Jawaban (C) karena `ratio of losses to wins` berarti perbandingan atau proportion antara kerugian dan kemenangan.\par
\endgroup
\par
\vspace{0.25\baselineskip}
\noindent\textbf{3.} Which of the following statements can be inferred from the first paragraph?
\begin{enumerate}[label=(\Alph*),leftmargin=*,itemsep=0.2em]
\item Temporary armies of farmers were not well trained.
\item Drafting farmers into armies was costly.
\item Heavy ground battles were won during combat.
\item Infantry was directed to the opposition for support.
\end{enumerate}
\par
\begingroup
\setlength{\parskip}{0pt}
\setlength{\parindent}{0pt}
\textbf{Review}\par
Answer: A\par
Confidence: medium\par
Evidence refs: lines 1-2, lines 7-8\par
Evidence quote: "landowners formed their own foot armies into which they drafted their tenants and hired hands ... this sporadic and untrained organization was ineffective"\par
Jawaban (A) karena pasukan awal direkrut dari tenants dan hired hands, lalu organisasi itu disebut `sporadic and untrained`, sehingga dapat diinferensikan bahwa tentara sementara dari rakyat biasa tidak terlatih baik.\par
Needs manual review: true\par
\endgroup
\par
\vspace{0.25\baselineskip}
\noindent\textbf{4.} In line 8, the word "sporadic" is closest in meaning to
\begin{enumerate}[label=(\Alph*),leftmargin=*,itemsep=0.2em]
\item spirited
\item splendid
\item irreverent
\item irregular
\end{enumerate}
\par
\begingroup
\setlength{\parskip}{0pt}
\setlength{\parindent}{0pt}
\textbf{Review}\par
Answer: D\par
Confidence: high\par
Evidence refs: line 8\par
Evidence quote: "Because this sporadic and untrained organization was ineffective"\par
Jawaban (D) karena `sporadic` dalam konteks ini berarti tidak tetap atau irregular.\par
\endgroup
\par
\vspace{0.25\baselineskip}
\noindent\textbf{5.} Which of the following statements about the Swiss pikemen is supported by the passage?
\begin{enumerate}[label=(\Alph*),leftmargin=*,itemsep=0.2em]
\item Their weapons and skills were ahead of their time.
\item Their gear ensured their fame as well-dressed soldiers.
\item The demand for their cavalry made them the best-paid army.
\item Their weapons were issued to nonprofessionals as well.
\end{enumerate}
\par
\begingroup
\setlength{\parskip}{0pt}
\setlength{\parindent}{0pt}
\textbf{Review}\par
Answer: A\par
Confidence: medium\par
Evidence refs: lines 10-13\par
Evidence quote: "The Swiss pikemen became the best-known mercenaries ... they practically invented a crude body armor ... Their tower shields proved indispensable ... their helmets progressed ..."\par
Jawaban (A) karena passage mendukung bahwa perlengkapan tempur Swiss pikemen maju dan penting pada masanya; opsi lain tidak didukung, meski kata `skills` tidak disebut secara langsung.\par
Needs manual review: true\par
\endgroup
\par
\vspace{0.25\baselineskip}
\noindent\textbf{6.} Where in the passage does the author state the reasons for the emergence of professional armies?
\begin{enumerate}[label=(\Alph*),leftmargin=*,itemsep=0.2em]
\item Lines 1-4
\item Lines 8-10
\item Lines 12-14
\item Lines 15-18
\end{enumerate}
\par
\begingroup
\setlength{\parskip}{0pt}
\setlength{\parindent}{0pt}
\textbf{Review}\par
Answer: B\par
Confidence: medium\par
Evidence refs: lines 7-10\par
Evidence quote: "Because this sporadic and untrained organization was ineffective, the ruling classes began to hire mercenaries ..."\par
Jawaban (B) karena alasan munculnya tentara profesional dijelaskan saat penulis menyebut organisasi lama tidak efektif lalu penguasa mulai mempekerjakan mercenaries; secara letak paling dekat dengan lines 8-10.\par
Needs manual review: true\par
\endgroup
\par
\vspace{0.25\baselineskip}
\noindent\textbf{7.} The author of the passage implies that the soldiers in mercenary armies were
\begin{enumerate}[label=(\Alph*),leftmargin=*,itemsep=0.2em]
\item not loyal
\item not effective
\item well guarded
\item well rounded
\end{enumerate}
\par
\begingroup
\setlength{\parskip}{0pt}
\setlength{\parindent}{0pt}
\textbf{Review}\par
Answer: A\par
Confidence: high\par
Evidence refs: lines 8-10\par
Evidence quote: "the troops sometimes deserted their employers if they could bank on a higher remuneration from the opposition"\par
Jawaban (A) karena passage menyatakan mereka bisa meninggalkan majikan demi bayaran lebih tinggi, jadi implikasinya mereka tidak loyal.\par
\endgroup
\par
\vspace{0.25\baselineskip}
\noindent\textbf{8.} In line 19, the phrase "these troops" refers to
\begin{enumerate}[label=(\Alph*),leftmargin=*,itemsep=0.2em]
\item the Swiss pikemen
\item military companies
\item almost all armies
\item Ottoman soldiers
\end{enumerate}
\par
\begingroup
\setlength{\parskip}{0pt}
\setlength{\parindent}{0pt}
\textbf{Review}\par
Answer: D\par
Confidence: high\par
Evidence refs: lines 18-20\par
Evidence quote: "the first full-size army of professional soldiers emerged in the Ottoman Empire. What set these troops apart ..."\par
Jawaban (D) karena `these troops` merujuk pada tentara profesional yang baru saja disebut muncul di Ottoman Empire, yaitu Ottoman soldiers.\par
\endgroup
\par
\vspace{0.25\baselineskip}
\noindent\textbf{9.} According to the passage, the first army of professionals was mobilized
\begin{enumerate}[label=(\Alph*),leftmargin=*,itemsep=0.2em]
\item only in peacetime
\item only in wartime
\item in times of anticipated war
\item both during war and during peace
\end{enumerate}
\par
\begingroup
\setlength{\parskip}{0pt}
\setlength{\parindent}{0pt}
\textbf{Review}\par
Answer: D\par
Confidence: high\par
Evidence refs: lines 14-16\par
Evidence quote: "the first full-size army of professional soldiers emerged ... these soldiers remained on duty in peacetime"\par
Jawaban (D) karena mereka disebut tetap bertugas saat masa damai, sehingga jelas tidak hanya dimobilisasi saat perang saja.\par
\endgroup
\par
\vspace{0.25\baselineskip}
\noindent\textbf{10.} In line 25, the word "undermined" is closest in meaning to
\begin{enumerate}[label=(\Alph*),leftmargin=*,itemsep=0.2em]
\item underestimated
\item reduced
\item undersized
\item shred
\end{enumerate}
\par
\begingroup
\setlength{\parskip}{0pt}
\setlength{\parindent}{0pt}
\textbf{Review}\par
Answer: B\par
Confidence: high\par
Evidence refs: lines 24-26\par
Evidence quote: "Cannons further widened the gap between the attacking and the defending lineups, and undermined the exclusivity of contact battles."\par
Jawaban (B) karena `undermined` di sini berarti mengurangi atau reduced the exclusivity of contact battles.\par
\endgroup
\par
\vspace{0.25\baselineskip}
\subsection*{Questions 11-18}
\begingroup
\small
\sloppy
\setlength{\parskip}{0pt}
\setlength{\parindent}{0pt}
Observations of nature gained a foothold in art in the 1860s and 1870s when painters in-\par
terested in science attempted to analyze the effects of light on color by means of physics. If\par
the goal of impressionist painters was to copy the visual qualities of sunlight at different an-\par
gles, they needed to reproduce light as it appears to the spectator when reflected from the\par
surfaces of structures. In painting, the effects of shade were conveyed by using small strokes\par
to minimize breaks between hues. The so-called divided color method appeared to grasp a\par
shimmering reflection of shadows when minimal portions of primary-color paints were ap-\par
plied directly to the canvas, instead of being blended on the palette.\par
Edouard Manet departed from the fairy-tale style of painting with its tacit symbolism\par
and centered his compositions around the visual reality of ordinary objects. Mary Cassatt fol-\par
lowed with her spontaneous and subtle portraits of children, and Edgar Degas depicted bal-\par
let dancers in their artful poses and the color schemes of their costumes in soft colors.\par
Postimpressionism built on the techniques developed by impressionists and supple-\par
mented it with keen insight into other dimensions of objects and scenes. Paul Gaugin chose\par
to disregard the classical conventions of composition, the application of color, and the shap-\par
ing of form and imitated primitivist art that upheld the beauty of native drawings in Tahiti.\par
Henri Matisse, created a unique style of poster graphics, deceptively simplistic in its rhythm\par
and texture. In his view, paintings were intended to brighten and improve reality, not copy\par
it. He noted that photography can accomplish this latter goal just as well, or even better.\par
\endgroup
\medskip
\noindent\textbf{11.} This passage probably comes from a longer work on
\begin{enumerate}[label=(\Alph*),leftmargin=*,itemsep=0.2em]
\item science and the fine arts
\item great masters of impressionist painting
\item light in the paintings of the 1800s
\item new techniques in the art of the 1800s
\end{enumerate}
\par
\begingroup
\setlength{\parskip}{0pt}
\setlength{\parindent}{0pt}
\textbf{Review}\par
Answer: D\par
Confidence: high\par
Evidence refs: lines 1-6, lines 11-17\par
Evidence quote: "The passage explains how painters analyzed light and color through physics, then describes impressionist and postimpressionist methods and styles."\par
Jawaban (D) karena isi passage membahas teknik baru dalam seni lukis abad 1800-an, terutama cara impresionis dan postimpresionis mengolah cahaya, warna, dan komposisi.\par
\endgroup
\par
\vspace{0.25\baselineskip}
\noindent\textbf{12.} The author of the passage implies that the goal of the impressionists was to
\begin{enumerate}[label=(\Alph*),leftmargin=*,itemsep=0.2em]
\item reproduce light and shadow exactly as they appeared
\item depict light as it appeared on different surfaces
\item copy the pattern of sun rays at different angles
\item reflect light and shadow from surfaces of structures
\end{enumerate}
\par
\begingroup
\setlength{\parskip}{0pt}
\setlength{\parindent}{0pt}
\textbf{Review}\par
Answer: B\par
Confidence: high\par
Evidence refs: lines 3-4\par
Evidence quote: "the goal of impressionist painters was to copy the visual qualities of sunlight at different angles ... reproduce light as it appears to the spectator when reflected from the surfaces of structures"\par
Jawaban (B) karena penulis menyatakan tujuan impresionis adalah mereproduksi cahaya sebagaimana terlihat oleh pengamat saat dipantulkan dari permukaan, jadi fokusnya pada tampilan cahaya di berbagai permukaan.\par
\endgroup
\par
\vspace{0.25\baselineskip}
\noindent\textbf{13.} In line 7, the word "shimmering" is closest in meaning to
\begin{enumerate}[label=(\Alph*),leftmargin=*,itemsep=0.2em]
\item gleaming
\item strong
\item trembling
\item spotless
\end{enumerate}
\par
\begingroup
\setlength{\parskip}{0pt}
\setlength{\parindent}{0pt}
\textbf{Review}\par
Answer: A\par
Confidence: high\par
Evidence refs: lines 6-8\par
Evidence quote: "The so-called divided color method appeared to grasp a shimmering reflection of shadows"\par
Jawaban (A) karena `shimmering` dalam konteks pantulan cahaya bayangan paling dekat maknanya dengan `gleaming` atau berkilau.\par
\endgroup
\par
\vspace{0.25\baselineskip}
\noindent\textbf{14.} In line 8, the word "blended" is closest in meaning to
\begin{enumerate}[label=(\Alph*),leftmargin=*,itemsep=0.2em]
\item replenished
\item placed
\item rendered
\item mixed
\end{enumerate}
\par
\begingroup
\setlength{\parskip}{0pt}
\setlength{\parindent}{0pt}
\textbf{Review}\par
Answer: D\par
Confidence: high\par
Evidence refs: lines 7-8\par
Evidence quote: "paints were applied directly to the canvas, instead of being blended on the palette"\par
Jawaban (D) karena `blended on the palette` berarti cat dicampur di palet terlebih dahulu; jadi `blended` paling dekat dengan `mixed`.\par
\endgroup
\par
\vspace{0.25\baselineskip}
\noindent\textbf{15.} What technique did painters employ to represent light as it appeared to the artist?
\begin{enumerate}[label=(\Alph*),leftmargin=*,itemsep=0.2em]
\item They used as little paint as was necessary.
\item They graded the shades of color in hues.
\item Their paint colors were dark and muted.
\item Their brush strokes were slow and cautious.
\end{enumerate}
\par
\begingroup
\setlength{\parskip}{0pt}
\setlength{\parindent}{0pt}
\textbf{Review}\par
Answer: A\par
Confidence: medium\par
Evidence refs: lines 9-10\par
Evidence quote: "minimal portions of primary-color paints were applied directly to the canvas, instead of being blended on the palette"\par
Jawaban (A) karena teknik yang dijelaskan adalah memakai porsi cat yang kecil secara langsung di kanvas. Pilihan ini paling dekat dengan evidence, meskipun wording soal dan opsi tidak sepenuhnya paralel dengan kalimat sumber.\par
Needs manual review: true\par
\endgroup
\par
\vspace{0.25\baselineskip}
\noindent\textbf{16.} It is implied in the passage that in the late 1800s artists painted
\begin{enumerate}[label=(\Alph*),leftmargin=*,itemsep=0.2em]
\item fantastic images and swirling action
\item luxurious ornaments and shapely figures
\item unusual textures and luscious colors
\item ordinary objects in varied intensities of light
\end{enumerate}
\par
\begingroup
\setlength{\parskip}{0pt}
\setlength{\parindent}{0pt}
\textbf{Review}\par
Answer: D\par
Confidence: medium\par
Evidence refs: lines 3-4, lines 7-10\par
Evidence quote: "copy the visual qualities of sunlight at different angles ... centered his compositions around the visual reality of ordinary objects"\par
Jawaban (D) karena passage menggabungkan dua ide: pelukis mengamati efek cahaya dari berbagai sudut dan Manet memusatkan lukisan pada realitas visual benda-benda biasa. Jadi inferensinya adalah objek biasa digambarkan dalam variasi intensitas cahaya.\par
Needs manual review: true\par
\endgroup
\par
\vspace{0.25\baselineskip}
\noindent\textbf{17.} It can be inferred from the passage that postimpressionism
\begin{enumerate}[label=(\Alph*),leftmargin=*,itemsep=0.2em]
\item developed into a trend parallel to impressionism
\item followed impressionism in the development of technique
\item flourished independently of impressionism
\item negated impressionist insight into objects and scenes
\end{enumerate}
\par
\begingroup
\setlength{\parskip}{0pt}
\setlength{\parindent}{0pt}
\textbf{Review}\par
Answer: B\par
Confidence: high\par
Evidence refs: lines 11-12\par
Evidence quote: "Postimpressionism built on the techniques developed by impressionists and supplemented it"\par
Jawaban (B) karena kalimat itu langsung menyebut postimpresionisme dibangun di atas teknik yang dikembangkan oleh impresionis, sehingga jelas datang setelah dan melanjutkan perkembangan teknik tersebut.\par
\endgroup
\par
\vspace{0.25\baselineskip}
\noindent\textbf{18.} According to the passage, Matisse saw the purpose of his art as
\begin{enumerate}[label=(\Alph*),leftmargin=*,itemsep=0.2em]
\item depicting life realistically
\item enhancing life and reality
\item showing life through photography
\item imparting rhythm to drawings
\end{enumerate}
\par
\begingroup
\setlength{\parskip}{0pt}
\setlength{\parindent}{0pt}
\textbf{Review}\par
Answer: B\par
Confidence: high\par
Evidence refs: lines 15-17\par
Evidence quote: "In his view, paintings were intended to brighten and improve reality, not copy it."\par
Jawaban (B) karena menurut passage Matisse memandang lukisan sebagai sarana untuk mencerahkan dan memperbaiki realitas, bukan sekadar menyalinnya secara realistis.\par
\endgroup
\par
\vspace{0.25\baselineskip}
\subsection*{Questions 19-30}
\begingroup
\small
\sloppy
\setlength{\parskip}{0pt}
\setlength{\parindent}{0pt}
The symptoms of hay fever include watery and itchy eyes and a runny, congested nose.\par
People suffering from hay fever may experience occasional wheezing and repeated bouts of\par
sneezing and may even lose their sense of smell. Some victims of hay fever may also have\par
stopped-up ears. About 30 percent of those who suffer from hay fever may develop the symp-\par
toms associated with periodic asthma or a sinus infection. The allergen-antibody theory\par
does not fully explain allergic reactions because the membranes and glands in eyes and ears\par
are controlled by the independent nervous system, which keeps these organs in balance. But\par
the independent nervous system itself is part of the emotional-response center and may\par
cause the feelings of anger, fear, resentment, and lack of self-confidence in reaction to al-\par
lergy-causing substances.\par
The most common cause of hay fever is the pollen of ragweed, which blossoms during the\par
summer and autumn. When airborne pollen particles, as well as mold, come into contact\par
with the victim's membranes, they can cause allergic reactions that release histamine and re-\par
sult in a virtual blockage of air passages. To prevent hay fever or to decrease the severity of its\par
symptoms, contact with the ragweed pollen should be reduced. Although some communi-\par
ties have attempted to eliminate the plants that cause the reactions, elimination programs\par
have not been successful because airborne pollen can travel considerable distances. Antihis-\par
tamine can help with short but severe attacks. Over extended periods of time, however, pa-\par
tients are prescribed a series of injections of the substance to which they are sensitive in order\par
to increase immunity and thus be relieved of the seasonal allergy\par
\endgroup
\medskip
\noindent\textbf{19.} It can be inferred from the passage that the phrase "hay fever" refers to
\begin{enumerate}[label=(\Alph*),leftmargin=*,itemsep=0.2em]
\item fodder for cattle
\item a seasonal discomfort
\item viral bacteria
\item a lung disease
\end{enumerate}
\par
\begingroup
\setlength{\parskip}{0pt}
\setlength{\parindent}{0pt}
\textbf{Review}\par
Answer: B\par
Confidence: high\par
Evidence refs: line 7, line 13\par
Evidence quote: "The most common cause of hay fever is the pollen of ragweed, which blossoms during the summer and autumn ... thus be relieved of the seasonal allergy"\par
Jawaban (B) karena passage menunjukkan hay fever berkaitan dengan musim tertentu dan bahkan disebut sebagai `seasonal allergy`, jadi maknanya adalah gangguan musiman.\par
\endgroup
\par
\vspace{0.25\baselineskip}
\noindent\textbf{20.} According to the passage, the symptoms of the allergy are predominantly
\begin{enumerate}[label=(\Alph*),leftmargin=*,itemsep=0.2em]
\item abdominal
\item intestinal
\item respiratory
\item chronic
\end{enumerate}
\par
\begingroup
\setlength{\parskip}{0pt}
\setlength{\parindent}{0pt}
\textbf{Review}\par
Answer: C\par
Confidence: high\par
Evidence refs: lines 1-3\par
Evidence quote: "The symptoms of hay fever include watery and itchy eyes and a runny, congested nose. People suffering from hay fever may experience occasional wheezing and repeated bouts of sneezing ... About 30 percent of those who suffer from hay fever may develop the symptoms associated with periodic asthma or a sinus infection."\par
Jawaban (C) karena gejala yang dominan adalah hidung tersumbat, wheezing, sneezing, asthma, dan sinus infection, yang terutama berkaitan dengan sistem pernapasan.\par
\endgroup
\par
\vspace{0.25\baselineskip}
\noindent\textbf{21.} What can be inferred from the first paragraph?
\begin{enumerate}[label=(\Alph*),leftmargin=*,itemsep=0.2em]
\item Hay fever may cause severe allergic reactions and even death.
\item The cause of allergic reactions has not been determined.
\item The nervous system balances allergic reactions.
\item People should not have an emotional response to allergic reactions.
\end{enumerate}
\par
\begingroup
\setlength{\parskip}{0pt}
\setlength{\parindent}{0pt}
\textbf{Review}\par
Answer: B\par
Confidence: medium\par
Evidence refs: lines 3-6\par
Evidence quote: "The allergen-antibody theory does not fully explain allergic reactions ... the independent nervous system itself is part of the emotional-response center and may cause the feelings of anger, fear, resentment, and lack of self-confidence in reaction to allergy-causing substances."\par
Jawaban (B) karena paragraf pertama menyatakan teori allergen-antibody `does not fully explain allergic reactions`, sehingga dapat disimpulkan penyebab atau mekanisme reaksi alergi belum sepenuhnya ditentukan.\par
Needs manual review: true\par
\endgroup
\par
\vspace{0.25\baselineskip}
\noindent\textbf{22.} According to the passage, patients suffering from hay fever may also experience
\begin{enumerate}[label=(\Alph*),leftmargin=*,itemsep=0.2em]
\item hunger pains
\item mood swings
\item nervous blockages
\item sensory perceptions
\end{enumerate}
\par
\begingroup
\setlength{\parskip}{0pt}
\setlength{\parindent}{0pt}
\textbf{Review}\par
Answer: B\par
Confidence: high\par
Evidence refs: lines 4-6\par
Evidence quote: "But the independent nervous system itself is part of the emotional-response center and may cause the feelings of anger, fear, resentment, and lack of self-confidence in reaction to allergy-causing substances."\par
Jawaban (B) karena passage menyebut efek emosional seperti anger, fear, resentment, dan lack of self-confidence, yang paling sesuai dengan mood swings.\par
\endgroup
\par
\vspace{0.25\baselineskip}
\noindent\textbf{23.} In line 9, the word "resentment" is closest in meaning to
\begin{enumerate}[label=(\Alph*),leftmargin=*,itemsep=0.2em]
\item reprieve
\item reprisal
\item acrimony
\item grief
\end{enumerate}
\par
\begingroup
\setlength{\parskip}{0pt}
\setlength{\parindent}{0pt}
\textbf{Review}\par
Answer: C\par
Confidence: high\par
Evidence refs: lines 9-10\par
Evidence quote: "cause the feelings of anger, fear, resentment, and lack of self-confidence in reaction to allergy-causing substances."\par
Jawaban (C) karena `resentment` paling dekat dengan rasa pahit atau kebencian, yaitu `acrimony`.\par
\endgroup
\par
\vspace{0.25\baselineskip}
\noindent\textbf{24.} It can be inferred from the passage that a frequent source of allergy-causing irritants can be
\begin{enumerate}[label=(\Alph*),leftmargin=*,itemsep=0.2em]
\item organic matter
\item larynx infections
\item human contact
\item ear membranes
\end{enumerate}
\par
\begingroup
\setlength{\parskip}{0pt}
\setlength{\parindent}{0pt}
\textbf{Review}\par
Answer: A\par
Confidence: medium\par
Evidence refs: lines 7-9\par
Evidence quote: "The most common cause of hay fever is the pollen of ragweed ... When airborne pollen particles, as well as mold, come into contact with the victim's membranes, they can cause allergic reactions"\par
Jawaban (A) karena sumber iritan yang disebut adalah pollen dan mold, keduanya merupakan bahan organik, sehingga `organic matter` adalah inferensi yang paling tepat.\par
Needs manual review: true\par
\endgroup
\par
\vspace{0.25\baselineskip}
\noindent\textbf{25.} According to the passage, the irritants are transported by
\begin{enumerate}[label=(\Alph*),leftmargin=*,itemsep=0.2em]
\item wind
\item food
\item travelers
\item air passages
\end{enumerate}
\par
\begingroup
\setlength{\parskip}{0pt}
\setlength{\parindent}{0pt}
\textbf{Review}\par
Answer: A\par
Confidence: high\par
Evidence refs: lines 7-11\par
Evidence quote: "When airborne pollen particles ... come into contact with the victim's membranes ... airborne pollen can travel considerable distances."\par
Jawaban (A) karena iritan disebut `airborne`, artinya dibawa melalui udara atau angin.\par
\endgroup
\par
\vspace{0.25\baselineskip}
\noindent\textbf{26.} In line 14, the word "blockage" is closest in meaning to
\begin{enumerate}[label=(\Alph*),leftmargin=*,itemsep=0.2em]
\item obstruction
\item bleeding
\item enlargement
\item dryness
\end{enumerate}
\par
\begingroup
\setlength{\parskip}{0pt}
\setlength{\parindent}{0pt}
\textbf{Review}\par
Answer: A\par
Confidence: high\par
Evidence refs: lines 13-14\par
Evidence quote: "release histamine and result in a virtual blockage of air passages."\par
Jawaban (A) karena `blockage` dalam konteks saluran udara berarti sumbatan atau `obstruction`.\par
\endgroup
\par
\vspace{0.25\baselineskip}
\noindent\textbf{27.} According to the passage, to avoid incidents of hay fever, patients need to
\begin{enumerate}[label=(\Alph*),leftmargin=*,itemsep=0.2em]
\item avoid interactions with other patients
\item avoid exposure to pollen
\item increase their self-confidence
\item take doses of prescribed medicine
\end{enumerate}
\par
\begingroup
\setlength{\parskip}{0pt}
\setlength{\parindent}{0pt}
\textbf{Review}\par
Answer: B\par
Confidence: high\par
Evidence refs: lines 9-11\par
Evidence quote: "To prevent hay fever or to decrease the severity of its symptoms, contact with the ragweed pollen should be reduced."\par
Jawaban (B) karena passage secara langsung mengatakan kontak dengan ragweed pollen harus dikurangi untuk mencegah hay fever.\par
\endgroup
\par
\vspace{0.25\baselineskip}
\noindent\textbf{28.} Which of the following is not mentioned in the passage as a cause of allergies?
\begin{enumerate}[label=(\Alph*),leftmargin=*,itemsep=0.2em]
\item pollen
\item mold
\item flowers
\item injections
\end{enumerate}
\par
\begingroup
\setlength{\parskip}{0pt}
\setlength{\parindent}{0pt}
\textbf{Review}\par
Answer: D\par
Confidence: high\par
Evidence refs: lines 7-9, lines 12-13\par
Evidence quote: "The most common cause of hay fever is the pollen of ragweed ... as well as mold ... patients are prescribed a series of injections ... in order to increase immunity"\par
Jawaban (D) karena pollen dan mold disebut sebagai penyebab alergi, sedangkan injections disebut sebagai pengobatan, bukan penyebab.\par
\endgroup
\par
\vspace{0.25\baselineskip}
\noindent\textbf{29.} It can be inferred from the passage that hay fever
\begin{enumerate}[label=(\Alph*),leftmargin=*,itemsep=0.2em]
\item has no effective antibodies
\item has no known cure
\item is rooted in the human psyche
\item can be likened to a breakdown
\end{enumerate}
\par
\begingroup
\setlength{\parskip}{0pt}
\setlength{\parindent}{0pt}
\textbf{Review}\par
Answer: B\par
Confidence: medium\par
Evidence refs: lines 9-13\par
Evidence quote: "To prevent hay fever or to decrease the severity of its symptoms ... Antihistamine can help with short but severe attacks ... patients are prescribed a series of injections ... to increase immunity and thus be relieved of the seasonal allergy"\par
Jawaban (B) karena passage hanya membahas pencegahan, pengurangan gejala, dan peredaan alergi musiman, bukan penyembuhan total; jadi inferensi terbaiknya adalah hay fever tidak memiliki cure yang diketahui dari passage ini.\par
Needs manual review: true\par
\endgroup
\par
\vspace{0.25\baselineskip}
\noindent\textbf{30.} A paragraph following this passage would most probably discuss
\begin{enumerate}[label=(\Alph*),leftmargin=*,itemsep=0.2em]
\item how the nervous system alerts patients
\item how the immune system reacts to allergens
\item what other diseases can be relieved by vaccines
\item what flowers are harmless to hay fever patients
\end{enumerate}
\par
\begingroup
\setlength{\parskip}{0pt}
\setlength{\parindent}{0pt}
\textbf{Review}\par
Answer: B\par
Confidence: medium\par
Evidence refs: lines 3-4, lines 8-13\par
Evidence quote: "The allergen-antibody theory does not fully explain allergic reactions ... release histamine ... patients are prescribed a series of injections ... to increase immunity"\par
Jawaban (B) karena passage sudah mengarah ke mekanisme alergi dan immunity, sehingga paragraf lanjutan paling mungkin membahas bagaimana sistem imun bereaksi terhadap allergens.\par
Needs manual review: true\par
\endgroup
\par
\vspace{0.25\baselineskip}
\subsection*{Questions 31-40}
\begingroup
\small
\sloppy
\setlength{\parskip}{0pt}
\setlength{\parindent}{0pt}
Prehistoric horses were far removed from the horses that Christopher Columbus brought\par
on his ships during his second voyage to the New World. Although fossil remains of "dawn\par
horses" have been excavated in several sites in Wyoming and New Mexico, these animals,\par
which were biologically different from contemporary horses, had become extinct millennia\par
before the onset of the Indian era. Although moviegoers visualize an Indian as a horse rider,\par
Indians were not familiar with horses until the Spanish brought them to Mexico, New Mex-\par
ico, Florida, and the West Indies in 1519. Those that escaped from the conquerors or were\par
left behind became the ancestors of the wild horses that still roam the southwestern regions\par
of the country. The Indian tribes scattered in the western plains began to breed horses about\par
1600.\par
The arrival of the horse produced a ripple effect throughout the Great Plains as the Indi-\par
ans living there were not nomadic and were engaged in rudimentary farming and graze-\par
land hunting. Tracking stampeding herds of buffalo and elk on foot was not the best way to\par
stock quantities of meat to adequately feed the entire tribe during the winter. However,\par
mounted on horses, the hunting teams could cover ground within a substantial distance\par
from their camps and transport their game back to be roasted, dried into jerky, or smoked\par
for preservation. The hunters responsible for tribe provisions stayed on the move almost con-\par
tinuously, replacing their earth-and-sod lodges with tepees. Horses carried not only their\par
riders but also their possessions and booty. The Blackfoot Indians of the Canadian plains\par
turned almost exclusively hunters, and the Crow split off from the mainstream Indian farm-\par
ing in favor of hunting. In fact, some of the Apache splinter groups abandoned agricultural\par
cultivation altogether.\par
The horse also drastically altered Indian warfare by allowing rapid maneuvering before,\par
during, and after skirmishes. With the advent of the horse, the Apache, Arapahoe, and\par
Cheyenne established themselves as a territorial monopoly in the Plains. Because Indians\par
did not have the wheel and had dragged their belongings from one settlement to another,\par
horses also enabled them to become more mobile and expedient during tribal migrations. In\par
fact, the Cheyenne abolished the custom of discarding belongings and tepee skins simply\par
because there were no means to transport them\par
\endgroup
\medskip
\noindent\textbf{31.} In line 3, the word "excavated" is closest in meaning to
\begin{enumerate}[label=(\Alph*),leftmargin=*,itemsep=0.2em]
\item exasperated
\item extinguished
\item hunted down
\item dug up
\end{enumerate}
\par
\begingroup
\setlength{\parskip}{0pt}
\setlength{\parindent}{0pt}
\textbf{Review}\par
Answer: D\par
Confidence: high\par
Evidence refs: line 3\par
Evidence quote: "Although fossil remains of dawn horses have been excavated in several sites in Wyoming and New Mexico"\par
Jawaban (D) karena `excavated` pada konteks fosil berarti digali dari dalam tanah, yaitu `dug up`.\par
\endgroup
\par
\vspace{0.25\baselineskip}
\noindent\textbf{32.} According to the passage, how many genetic species of horses are known today?
\begin{enumerate}[label=(\Alph*),leftmargin=*,itemsep=0.2em]
\item One
\item Two
\item Three
\item Four
\end{enumerate}
\par
\begingroup
\setlength{\parskip}{0pt}
\setlength{\parindent}{0pt}
\textbf{Review}\par
Answer: B\par
Confidence: medium\par
Evidence refs: lines 1-2\par
Evidence quote: "Prehistoric horses were far removed from the horses that Christopher Columbus brought ... these animals, which were biologically different from contemporary horses, had become extinct"\par
Jawaban (B) karena passage menyebut dua kelompok yang berbeda secara biologis: `dawn horses` prasejarah dan kuda kontemporer. Jadi, dari passage dapat diinferensikan ada dua species/jenis biologis yang diketahui.\par
Needs manual review: true\par
\endgroup
\par
\vspace{0.25\baselineskip}
\noindent\textbf{33.} In line 7, the word "those" refers to
\begin{enumerate}[label=(\Alph*),leftmargin=*,itemsep=0.2em]
\item West Indies
\item The Spanish
\item horses,
\item Indians
\end{enumerate}
\par
\begingroup
\setlength{\parskip}{0pt}
\setlength{\parindent}{0pt}
\textbf{Review}\par
Answer: C\par
Confidence: high\par
Evidence refs: lines 6-8\par
Evidence quote: "Indians were not familiar with horses until the Spanish brought them ... Those that escaped from the conquerors or were left behind became the ancestors of the wild horses"\par
Jawaban (C) karena kata `Those` merujuk pada `horses` yang dibawa oleh orang Spanyol lalu ada yang lepas atau tertinggal.\par
\endgroup
\par
\vspace{0.25\baselineskip}
\noindent\textbf{34.} According to the passage, American Indians
\begin{enumerate}[label=(\Alph*),leftmargin=*,itemsep=0.2em]
\item tamed horses in the early 1500s
\item farmed with horses in the 1500s
\item were exposed to horses in the 1500s
\item have ridden horses since prehistoric times
\end{enumerate}
\par
\begingroup
\setlength{\parskip}{0pt}
\setlength{\parindent}{0pt}
\textbf{Review}\par
Answer: C\par
Confidence: high\par
Evidence refs: lines 2-5\par
Evidence quote: "Indians were not familiar with horses until the Spanish brought them to Mexico, New Mexico, Florida, and the West Indies in 1519"\par
Jawaban (C) karena teks menyatakan Indian baru mengenal kuda ketika bangsa Spanyol membawanya pada 1519, yaitu abad ke-16.\par
\endgroup
\par
\vspace{0.25\baselineskip}
\noindent\textbf{35.} The author of the passage probably believes that the popular image of American Indians before the arrival of Europeans
\begin{enumerate}[label=(\Alph*),leftmargin=*,itemsep=0.2em]
\item is not theoretically viable
\item cannot be realistically described
\item cannot be discussed briefly
\item is not historically accurate
\end{enumerate}
\par
\begingroup
\setlength{\parskip}{0pt}
\setlength{\parindent}{0pt}
\textbf{Review}\par
Answer: D\par
Confidence: high\par
Evidence refs: lines 2-3\par
Evidence quote: "Although moviegoers visualize an Indian as a horse rider, Indians were not familiar with horses until the Spanish brought them"\par
Jawaban (D) karena penulis secara jelas mengoreksi gambaran populer bahwa Indian selalu menunggang kuda; secara historis itu tidak akurat sebelum kedatangan orang Eropa.\par
\endgroup
\par
\vspace{0.25\baselineskip}
\noindent\textbf{36.} According to the passage, after the arrival of Europeans, the Indian tribes inhabiting the Great Plains
\begin{enumerate}[label=(\Alph*),leftmargin=*,itemsep=0.2em]
\item herded undomesticated buffalo
\item played complicated hunting games
\item had sedentary and tranquil life-styles
\item improved their hunting techniques
\end{enumerate}
\par
\begingroup
\setlength{\parskip}{0pt}
\setlength{\parindent}{0pt}
\textbf{Review}\par
Answer: D\par
Confidence: high\par
Evidence refs: lines 6-10\par
Evidence quote: "mounted on horses, the hunting teams could cover ground within a substantial distance from their camps and transport their game back"\par
Jawaban (D) karena setelah ada kuda, tim pemburu bisa menempuh jarak lebih jauh dan membawa hasil buruan kembali, sehingga teknik berburu mereka menjadi lebih efektif.\par
\endgroup
\par
\vspace{0.25\baselineskip}
\noindent\textbf{37.} In line 17, the word "provisions" is closest in meaning to
\begin{enumerate}[label=(\Alph*),leftmargin=*,itemsep=0.2em]
\item supplies
\item health
\item weapons
\item attire
\end{enumerate}
\par
\begingroup
\setlength{\parskip}{0pt}
\setlength{\parindent}{0pt}
\textbf{Review}\par
Answer: A\par
Confidence: high\par
Evidence refs: lines 16-17\par
Evidence quote: "The hunters responsible for tribe provisions stayed on the move almost continuously"\par
Jawaban (A) karena `provisions` dalam konteks hasil buruan untuk memberi makan suku berarti persediaan atau `supplies`.\par
\endgroup
\par
\vspace{0.25\baselineskip}
\noindent\textbf{38.} According to the passage, American Indians invented various methods for
\begin{enumerate}[label=(\Alph*),leftmargin=*,itemsep=0.2em]
\item dislocating their traps
\item communicating over great distances
\item conducting their hostile excursions
\item keeping their possessions
\end{enumerate}
\par
\begingroup
\setlength{\parskip}{0pt}
\setlength{\parindent}{0pt}
\textbf{Review}\par
Answer: D\par
Confidence: medium\par
Evidence refs: lines 12-18\par
Evidence quote: "Horses carried not only their riders but also their possessions ... horses also enabled them to become more mobile ... the Cheyenne abolished the custom of discarding belongings"\par
Jawaban (D) karena bagian ini menekankan bahwa dengan kuda mereka dapat membawa dan tidak lagi membuang barang-barang milik mereka. Pilihan ini paling sesuai meskipun redaksi soal kurang langsung terhadap teks.\par
Needs manual review: true\par
\endgroup
\par
\vspace{0.25\baselineskip}
\noindent\textbf{39.} It can be inferred from the passage that Indians did NOT
\begin{enumerate}[label=(\Alph*),leftmargin=*,itemsep=0.2em]
\item accrue tribal wealth
\item assign sustenance tasks
\item pursue stampedes
\item use covered wagons
\end{enumerate}
\par
\begingroup
\setlength{\parskip}{0pt}
\setlength{\parindent}{0pt}
\textbf{Review}\par
Answer: D\par
Confidence: high\par
Evidence refs: line 18\par
Evidence quote: "Because Indians did not have the wheel"\par
Jawaban (D) karena jika mereka tidak memiliki roda, maka dapat disimpulkan mereka tidak menggunakan kendaraan beroda seperti covered wagons.\par
\endgroup
\par
\vspace{0.25\baselineskip}
\noindent\textbf{40.} It can be inferred from the passage that the arrival of horses in the Americas
\begin{enumerate}[label=(\Alph*),leftmargin=*,itemsep=0.2em]
\item led to the dispersal of Indian tribes throughout the continent
\item made Indian tribes relinquish their territorial monopolies
\item altered the future course of the Indian way of life
\item shattered the advancement of the Indian culture
\end{enumerate}
\par
\begingroup
\setlength{\parskip}{0pt}
\setlength{\parindent}{0pt}
\textbf{Review}\par
Answer: C\par
Confidence: medium\par
Evidence refs: lines 6-18\par
Evidence quote: "The arrival of the horse produced a ripple effect throughout the Great Plains ... The horse also drastically altered Indian warfare ... horses also enabled them to become more mobile"\par
Jawaban (C) karena teks berulang kali menunjukkan bahwa kedatangan kuda mengubah perburuan, pola hidup, peperangan, dan mobilitas suku Indian; jadi arah kehidupan mereka berubah secara besar. Ini merupakan inferensi menyeluruh dari passage.\par
Needs manual review: true\par
\endgroup
\par
\vspace{0.25\baselineskip}
\subsection*{Questions 41-50}
\begingroup
\small
\sloppy
\setlength{\parskip}{0pt}
\setlength{\parindent}{0pt}
Lighthouses and lightships employ signal lights and foghorns to warn boaters of shoals\par
and of oncoming foul weather. In recreational marinas, Coast Guard stations and yacht\par
clubs inform boaters of weather and water conditions with storm flag signals displayed dur-\par
ing the daytime. Amateur boaters are required to acquaint themselves with the signals that\par
can make them aware of oncoming storms and small craft advisories. Boaters who can recog-\par
nize an approaching storm when no warnings are posted have the advantage of time when\par
heading for shelter. Compasses and marine charts identify the locations of anchored whistles\par
and floating buoys that have battery-powered lights, bells, or horns to sound or show warn-\par
ings with the onset of high winds. On the marine charts, radio beacons and flashing buoys\par
are assigned numbers that help navigators to identify their locations and to mark the edges\par
of channels, underwater obstructions, reefs, and wrecks.\par
In each section of the coastline, lighthouses emit characteristic lights, published in light\par
lists, so that mariners can zero in on their bearings by observing the beam pattern and con-\par
sulting the list. "Making lights" signal approaching vessels to make land, and "leading\par
lights" guide navigators into bays and harbors along navigable waterways.\par
The lighthouse tower, constructed on solid rock and pneumatic caissons, contains light-\par
ing mechanisms, engines, and spare parts, as well as the keeper's quarters. Saucerlike Fres-\par
nel lenses with pentagonal prisms project light at irregular, alternating intervals, while\par
sealed-beam lenses rotate at varying speeds, similarly to the traditional long-range search\par
light. The older, barrel-shaped lenses and kerosene-burning lamps, made up of prisms and\par
glass panels attached to a metal frame, have been replaced by acetylene gas burners and in-\par
candescent lamps that can be operated either manually or automatically and that shut off at\par
daybreak\par
\endgroup
\medskip
\noindent\textbf{41.} What is the best title for the passage?
\begin{enumerate}[label=(\Alph*),leftmargin=*,itemsep=0.2em]
\item The Function of Naval Signals and Buoys
\item Marine Installations for Boaters' Safety
\item Lighthouses and Maritime Signaling Devices
\item Lighthouses and Occulting Instrumentation
\end{enumerate}
\par
\begingroup
\setlength{\parskip}{0pt}
\setlength{\parindent}{0pt}
\textbf{Review}\par
Answer: C\par
Confidence: high\par
Evidence refs: lines 1-17\par
Evidence quote: "Passage membahas lighthouses and lightships, lalu floating buoys, radio beacons, dan bagaimana lighthouses emit characteristic lights untuk membantu navigasi."\par
Jawaban (C) karena isi passage berfokus pada mercusuar sekaligus perangkat pensinyalan maritim lain seperti lightships, buoys, beacons, dan signal lights; itu paling sesuai dengan judul umum passage.\par
\endgroup
\par
\vspace{0.25\baselineskip}
\noindent\textbf{42.} According to the passage, what is the purpose of lighthouses?
\begin{enumerate}[label=(\Alph*),leftmargin=*,itemsep=0.2em]
\item To measure weather and water conditions
\item To lead navigators to their destination
\item To flag down oncoming vessels
\item To display light-coded messages
\end{enumerate}
\par
\begingroup
\setlength{\parskip}{0pt}
\setlength{\parindent}{0pt}
\textbf{Review}\par
Answer: B\par
Confidence: high\par
Evidence refs: lines 9-12\par
Evidence quote: "lighthouses emit characteristic lights ... so that mariners can zero in on their bearings and leading lights guide navigators into bays and harbors."\par
Jawaban (B) karena passage menyebut mercusuar membantu mariners menentukan bearing dan memandu navigators masuk ke teluk serta pelabuhan.\par
\endgroup
\par
\vspace{0.25\baselineskip}
\noindent\textbf{43.} According to the passage, buoys can be best described as
\begin{enumerate}[label=(\Alph*),leftmargin=*,itemsep=0.2em]
\item hanging constructions
\item electrical horns
\item floating devices
\item underwater anchors
\end{enumerate}
\par
\begingroup
\setlength{\parskip}{0pt}
\setlength{\parindent}{0pt}
\textbf{Review}\par
Answer: C\par
Confidence: high\par
Evidence refs: lines 5-8\par
Evidence quote: "floating buoys that have battery-powered lights, bells, or horns"\par
Jawaban (C) karena passage secara eksplisit menyebut buoys sebagai `floating buoys`, jadi buoys adalah perangkat yang mengapung.\par
\endgroup
\par
\vspace{0.25\baselineskip}
\noindent\textbf{44.} Why do recreational mariners need to be familiar with the Coast Guard signaling system?
\begin{enumerate}[label=(\Alph*),leftmargin=*,itemsep=0.2em]
\item To learn navigation and rowing
\item To chart the course of storms
\item To learn about oncoming storms
\item To overtake rivals in boating races
\end{enumerate}
\par
\begingroup
\setlength{\parskip}{0pt}
\setlength{\parindent}{0pt}
\textbf{Review}\par
Answer: C\par
Confidence: high\par
Evidence refs: lines 2-5\par
Evidence quote: "Amateur boaters are required to acquaint themselves with the signals that can make them aware of oncoming storms and small craft advisories."\par
Jawaban (C) karena recreational mariners perlu memahami sistem sinyal agar tahu adanya badai yang mendekat dan small craft advisories.\par
\endgroup
\par
\vspace{0.25\baselineskip}
\noindent\textbf{45.} In line 6, the phrase "the advantage of time" means most nearly
\begin{enumerate}[label=(\Alph*),leftmargin=*,itemsep=0.2em]
\item an early warning
\item a stable bearing
\item a timed landing
\item a timely arrival
\end{enumerate}
\par
\begingroup
\setlength{\parskip}{0pt}
\setlength{\parindent}{0pt}
\textbf{Review}\par
Answer: A\par
Confidence: high\par
Evidence refs: lines 5-7\par
Evidence quote: "Boaters who can recognize an approaching storm when no warnings are posted have the advantage of time when heading for shelter."\par
Jawaban (A) karena `the advantage of time` di konteks ini berarti mendapat peringatan lebih awal, sehingga ada waktu untuk menuju tempat berlindung.\par
\endgroup
\par
\vspace{0.25\baselineskip}
\noindent\textbf{46.} The author of the passage implies that flashing lights from lighthouses function as
\begin{enumerate}[label=(\Alph*),leftmargin=*,itemsep=0.2em]
\item a characteristic marine elevation
\item beams with a constant rotating speed
\item an indication of the lens thickness
\item an identification of its location
\end{enumerate}
\par
\begingroup
\setlength{\parskip}{0pt}
\setlength{\parindent}{0pt}
\textbf{Review}\par
Answer: D\par
Confidence: high\par
Evidence refs: lines 9-12\par
Evidence quote: "lighthouses emit characteristic lights, published in light lists, so that mariners can zero in on their bearings by observing the beam pattern"\par
Jawaban (D) karena pola cahaya yang khas dipakai mariners untuk mengenali bearing dan lokasi mercusuar tersebut.\par
\endgroup
\par
\vspace{0.25\baselineskip}
\noindent\textbf{47.} According to the passage, lighthouses can assist navigators in
\begin{enumerate}[label=(\Alph*),leftmargin=*,itemsep=0.2em]
\item identifying underwater obstructions
\item preparing their boats for advisories.
\item finding their positions in the open sea
\item ignoring their charts of inland channels
\end{enumerate}
\par
\begingroup
\setlength{\parskip}{0pt}
\setlength{\parindent}{0pt}
\textbf{Review}\par
Answer: C\par
Confidence: high\par
Evidence refs: lines 9-12\par
Evidence quote: "mariners can zero in on their bearings by observing the beam pattern"\par
Jawaban (C) karena passage menjelaskan bahwa mercusuar membantu mariners menentukan bearing mereka, yaitu menemukan posisi mereka saat bernavigasi.\par
\endgroup
\par
\vspace{0.25\baselineskip}
\noindent\textbf{48.} In line 13, the phrase "zero in on" is closest in meaning to
\begin{enumerate}[label=(\Alph*),leftmargin=*,itemsep=0.2em]
\item cast about for
\item disguise
\item pinpoint
\item point out
\end{enumerate}
\par
\begingroup
\setlength{\parskip}{0pt}
\setlength{\parindent}{0pt}
\textbf{Review}\par
Answer: C\par
Confidence: high\par
Evidence refs: line 13\par
Evidence quote: "mariners can zero in on their bearings by observing the beam pattern"\par
Jawaban (C) karena `zero in on` di konteks ini berarti menentukan secara tepat atau pinpoint posisi atau bearing mereka.\par
\endgroup
\par
\vspace{0.25\baselineskip}
\noindent\textbf{49.} In line 15, the word "navigable" is closest in meaning to
\begin{enumerate}[label=(\Alph*),leftmargin=*,itemsep=0.2em]
\item nauseating
\item passable
\item pernicious
\item calamitous
\end{enumerate}
\par
\begingroup
\setlength{\parskip}{0pt}
\setlength{\parindent}{0pt}
\textbf{Review}\par
Answer: B\par
Confidence: high\par
Evidence refs: lines 14-15\par
Evidence quote: "guide navigators into bays and harbors along navigable waterways"\par
Jawaban (B) karena `navigable waterways` berarti jalur air yang dapat dilalui kapal, jadi maknanya paling dekat adalah passable.\par
\endgroup
\par
\vspace{0.25\baselineskip}
\noindent\textbf{50.} What can be inferred from the last paragraph?
\begin{enumerate}[label=(\Alph*),leftmargin=*,itemsep=0.2em]
\item Lighthouses do not contain a living space for maintenance personnel.
\item Lighthouse beams are projected intermittently during nighttime.
\item Technological expertise is expected of the lighthouse maintenance personnel.
\item Lighthouse technology is outdated and should have been replaced.
\end{enumerate}
\par
\begingroup
\setlength{\parskip}{0pt}
\setlength{\parindent}{0pt}
\textbf{Review}\par
Answer: B\par
Confidence: high\par
Evidence refs: lines 13-17\par
Evidence quote: "Fresnel lenses ... project light at irregular, alternating intervals ... lamps can be operated either manually or automatically and shut off at daybreak."\par
Jawaban (B) karena paragraf terakhir menjelaskan bahwa berkas cahaya diproyeksikan pada interval yang tidak tetap atau alternating, sehingga dapat diinferensikan bahwa sinar mercusuar tampak berkedip/intermittent pada malam hari.\par
\endgroup
\par
\vspace{0.25\baselineskip}
\newpage
\section*{Answer Key Appendix}
\subsection*{Listening}
\begin{itemize}[leftmargin=*]
\item 1: C
\item 2: C
\item 3: B
\item 4: A
\item 5: B
\item 6: C
\item 7: B
\item 8: A
\item 9: C
\item 10: B
\item 11: C
\item 12: B
\item 13: B
\item 14: D
\item 15: B
\item 16: B
\item 17: D
\item 18: C
\item 19: C
\item 20: B
\item 21: B
\item 22: C
\item 23: C
\item 24: B
\item 25: A
\item 26: B
\item 27: D
\item 28: C
\item 29: A
\item 30: C
\item 31: C
\item 32: C
\item 33: B
\item 34: B
\item 35: C
\item 36: B
\item 37: A
\item 38: B
\item 39: D
\item 40: D
\item 41: A
\item 42: C
\item 43: A
\item 44: B
\item 45: C
\item 46: D
\item 47: B
\item 48: A
\item 49: C
\item 50: D
\end{itemize}
\subsection*{Structure}
\begin{itemize}[leftmargin=*]
\item 1: B
\item 2: B
\item 3: C
\item 4: C
\item 5: C
\item 6: B
\item 7: D
\item 8: B
\item 9: B
\item 10: C
\item 11: D
\item 12: B
\item 13: C
\item 14: B
\item 15: D
\item 16: D
\item 17: D
\item 18: C
\item 19: B
\item 20: D
\item 21: A
\item 22: A
\item 23: A
\item 24: C
\item 25: D
\item 26: B
\item 27: C
\item 28: A
\item 29: D
\item 30: C
\item 31: D
\item 32: A
\item 33: B
\item 34: B
\item 35: C
\item 36: D
\item 37: D
\item 38: C
\item 39: C
\item 40: B
\end{itemize}
\subsection*{Reading}
\begin{itemize}[leftmargin=*]
\item 1: D
\item 2: C
\item 3: A
\item 4: D
\item 5: A
\item 6: B
\item 7: A
\item 8: D
\item 9: D
\item 10: B
\item 11: D
\item 12: B
\item 13: A
\item 14: D
\item 15: A
\item 16: D
\item 17: B
\item 18: B
\item 19: B
\item 20: C
\item 21: B
\item 22: B
\item 23: C
\item 24: A
\item 25: A
\item 26: A
\item 27: B
\item 28: D
\item 29: B
\item 30: B
\item 31: D
\item 32: B
\item 33: C
\item 34: C
\item 35: D
\item 36: D
\item 37: A
\item 38: D
\item 39: D
\item 40: C
\item 41: C
\item 42: B
\item 43: C
\item 44: C
\item 45: A
\item 46: D
\item 47: C
\item 48: C
\item 49: B
\item 50: B
\end{itemize}
\end{document}
